% Options for packages loaded elsewhere
\PassOptionsToPackage{unicode}{hyperref}
\PassOptionsToPackage{hyphens}{url}
\PassOptionsToPackage{dvipsnames,svgnames,x11names}{xcolor}
%
\documentclass[
  10pt,
]{article}
\usepackage{amsmath,amssymb}
\usepackage{iftex}
\ifPDFTeX
  \usepackage[T1]{fontenc}
  \usepackage[utf8]{inputenc}
  \usepackage{textcomp} % provide euro and other symbols
\else % if luatex or xetex
  \usepackage{unicode-math} % this also loads fontspec
  \defaultfontfeatures{Scale=MatchLowercase}
  \defaultfontfeatures[\rmfamily]{Ligatures=TeX,Scale=1}
\fi
\usepackage{lmodern}
\ifPDFTeX\else
  % xetex/luatex font selection
\fi
% Use upquote if available, for straight quotes in verbatim environments
\IfFileExists{upquote.sty}{\usepackage{upquote}}{}
\IfFileExists{microtype.sty}{% use microtype if available
  \usepackage[]{microtype}
  \UseMicrotypeSet[protrusion]{basicmath} % disable protrusion for tt fonts
}{}
\makeatletter
\@ifundefined{KOMAClassName}{% if non-KOMA class
  \IfFileExists{parskip.sty}{%
    \usepackage{parskip}
  }{% else
    \setlength{\parindent}{0pt}
    \setlength{\parskip}{6pt plus 2pt minus 1pt}}
}{% if KOMA class
  \KOMAoptions{parskip=half}}
\makeatother
\usepackage{xcolor}
\usepackage[margin=2.3cm]{geometry}
\usepackage{longtable,booktabs,array}
\usepackage{calc} % for calculating minipage widths
% Correct order of tables after \paragraph or \subparagraph
\usepackage{etoolbox}
\makeatletter
\patchcmd\longtable{\par}{\if@noskipsec\mbox{}\fi\par}{}{}
\makeatother
% Allow footnotes in longtable head/foot
\IfFileExists{footnotehyper.sty}{\usepackage{footnotehyper}}{\usepackage{footnote}}
\makesavenoteenv{longtable}
\setlength{\emergencystretch}{3em} % prevent overfull lines
\providecommand{\tightlist}{%
  \setlength{\itemsep}{0pt}\setlength{\parskip}{0pt}}
\setcounter{secnumdepth}{-\maxdimen} % remove section numbering
% Extra preamble for the pandoc LaTeX template (pdflatex).
\usepackage{amsmath,amssymb}
\usepackage{graphicx}
\usepackage{cite}                 % numeric citations, ranges compressed: [3--5]
\usepackage{caption}
\captionsetup{font=small,labelsep=period,justification=raggedright,singlelinecheck=false}
\usepackage{etoolbox}
\usepackage{framed}
\definecolor{shadecolor}{gray}{0.95}
\AtBeginEnvironment{longtable}{\small}
\setlength{\emergencystretch}{3em}
% allow line breaks after underscores (column names in \texttt)
\DeclareRobustCommand{\_}{\textunderscore\hspace{0pt}}
% characters used in the text that the utf8 input encoding does not define
\DeclareUnicodeCharacter{2212}{\ensuremath{-}}
\DeclareUnicodeCharacter{2192}{\ensuremath{\rightarrow}}
\DeclareUnicodeCharacter{03C1}{\ensuremath{\rho}}
\DeclareUnicodeCharacter{03B2}{\ensuremath{\beta}}
\DeclareUnicodeCharacter{0394}{\ensuremath{\Delta}}
\DeclareUnicodeCharacter{03C3}{\ensuremath{\sigma}}
\DeclareUnicodeCharacter{03BA}{\ensuremath{\kappa}}
\DeclareUnicodeCharacter{2265}{\ensuremath{\geq}}
\DeclareUnicodeCharacter{2264}{\ensuremath{\leq}}
\DeclareUnicodeCharacter{2248}{\ensuremath{\approx}}
\DeclareUnicodeCharacter{00B3}{\ensuremath{{}^{3}}}
\DeclareUnicodeCharacter{2011}{\mbox{-}}
\DeclareUnicodeCharacter{202F}{\,}
\ifLuaTeX
  \usepackage{selnolig}  % disable illegal ligatures
\fi
\IfFileExists{bookmark.sty}{\usepackage{bookmark}}{\usepackage{hyperref}}
\IfFileExists{xurl.sty}{\usepackage{xurl}}{} % add URL line breaks if available
\urlstyle{same}
\hypersetup{
  pdftitle={{[}Title{]}},
  colorlinks=true,
  linkcolor={blue},
  filecolor={Maroon},
  citecolor={Blue},
  urlcolor={blue},
  pdfcreator={LaTeX via pandoc}}

\title{{[}Title{]}}
\author{}
\date{}

\begin{document}
\maketitle

\nocite{alford2017,
  chaudhury2010,
  lee1971,
  mitternacht2016,
  tien2013,
  olsson2011,
  sondergaard2011,
  chai2024,
  evans2021,
  jumper2021,
  ahern2026,
  burley2025,
  berman2000,
  uniprot2023,
  dana2019,
  ribeiro2018,
  anishchenko2025,
  hanley1982,
  mann1947,
  efron1979,
  field2007,
  westfall1993,
  kim2026,
  carlin2016,
  carlin2017,
  henikoff1992,
  spearman1904}

\section*{Abstract}
\addcontentsline{toc}{section}{Abstract}

\section{1
Introduction}

\section{2 Methods}

\subsection{2.1
Experiment overview}

Four tests, one in
each of 2.2 to 2.5,
ask whether the
in-silico metrics
used to filter
designed enzymes
respond to catalysis
or to something
simpler, such as how
close a change is to
the active site:

\begin{itemize}
\tightlist
\item
  \textbf{Activity
  discrimination
  test (2.2; results
  in 3.1).} The
  metrics are ranked
  by how well they
  separate real dead
  from active
  proteins, 21
  zymogen--mature
  enzyme pairs and
  192 plated de novo
  designs of which
  16 are active.
\item
  \textbf{Substrate
  discrimination
  test (2.3; results
  in 3.2).} A
  structure
  predictor is given
  each of 55 natural
  enzymes with its
  own substrate and
  with a wrong
  ligand, and the
  prediction-based
  metrics are ranked
  by how well they
  separate the two
  predictions.
\item
  \textbf{Activity
  ranking test (2.4;
  results in 3.3).}
  The metrics are
  ranked by how well
  their order of 432
  BglB variants and
  of 16 plated
  designs agrees
  with the measured
  activity, against
  declared trivial
  baselines.
\item
  \textbf{Detection
  ability test (2.5;
  results in 3.4).}
  Catalytic residues
  are deliberately
  damaged in up to
  195 natural
  enzymes, and the
  metrics are ranked
  by whether they
  move more than
  when the same
  damage is made at
  a matched
  non-catalytic site
  of the same
  protein.
\end{itemize}

\textbf{Metrics.}
The same panel is
scored in every
test: 29
structure-space
metrics, which read
atomic coordinates
and are used in 2.2,
2.4 and 2.5, and 4
prediction-based
metrics, which read
the output of a
structure predictor
that is given a
sequence and a
ligand and are used
in 2.2 to 2.5 (Table
1). Numbers in
square brackets are
references (see
References); tables
with an S number are
supplementary
tables, described in
the Supplementary
PDF and released as
CSV files
(\texttt{main/} and
\texttt{supplementary/}).

\begin{longtable}[]{@{}
  >{\raggedright\arraybackslash}p{(\columnwidth - 6\tabcolsep) * \real{0.1613}}
  >{\raggedright\arraybackslash}p{(\columnwidth - 6\tabcolsep) * \real{0.3226}}
  >{\raggedright\arraybackslash}p{(\columnwidth - 6\tabcolsep) * \real{0.3548}}
  >{\raggedright\arraybackslash}p{(\columnwidth - 6\tabcolsep) * \real{0.1613}}@{}}
\caption*{\textbf{Table
1. The 33 reported
metrics and the six
reference rows.} In
a structure,
``perturbed
residues'' are the
residues that were
changed; in the
control arm they are
the control
residues, so a
site-scoped metric
never knows which
residues are
catalytic --- it
only knows which
residues it was
handed. In 2.2 and
2.4 the scored
residues are the
catalytic residues.
The third column
says where each
definition comes
from.}\tabularnewline
\toprule\noalign{}
\begin{minipage}[b]{\linewidth}\raggedright
Group
\end{minipage} &
\begin{minipage}[b]{\linewidth}\raggedright
Metrics
\end{minipage} &
\begin{minipage}[b]{\linewidth}\raggedright
Reads, and where it
is defined
\end{minipage} &
\begin{minipage}[b]{\linewidth}\raggedright
Scored over
\end{minipage} \\
\midrule\noalign{}
\endfirsthead
\toprule\noalign{}
\begin{minipage}[b]{\linewidth}\raggedright
Group
\end{minipage} &
\begin{minipage}[b]{\linewidth}\raggedright
Metrics
\end{minipage} &
\begin{minipage}[b]{\linewidth}\raggedright
Reads, and where it
is defined
\end{minipage} &
\begin{minipage}[b]{\linewidth}\raggedright
Scored over
\end{minipage} \\
\midrule\noalign{}
\endhead
\bottomrule\noalign{}
\endlastfoot
Rosetta energy,
site-scoped (15) &
attraction,
repulsion,
intra-residue
repulsion,
solvation, lk-ball
solvation,
electrostatics,
hydrogen bonds
(backbone--side-chain,
side-chain--side-chain),
rotamer (Dunbrack),
Ramachandran,
amino-acid
propensity, omega
torsion, net van der
Waals (repulsion
plus attraction),
site total, worst
single residue &
coordinates; the
per-residue terms of
the REF2015 energy
function
\cite{alford2017}
computed with
PyRosetta
\cite{chaudhury2010}
and summed over the
scored residues
(site total: the sum
of the residue
totals; worst
residue: the largest
residue total) & the
perturbed
residues \\
Catalytic geometry
(8) & fraction
buried, relative
solvent-accessible
area (mean, max),
total
solvent-accessible
area, pairwise
distance (mean,
max), radius of
gyration, polar
contacts &
coordinates; SASA by
the Lee--Richards
algorithm
\cite{lee1971} in
FreeSASA
\cite{mitternacht2016}
(probe 1.4 Å, waters
excluded), relative
to the maximum
accessibility of the
residue type
\cite{tien2013};
fraction buried is
the share of the
scored residues with
relative SASA below
0.10; distances and
radius of gyration
are over all pairs,
or all atoms, of the
side-chain heavy
atoms of the scored
residues (backbone
atoms for Gly); a
polar contact is a
pair of polar
side-chain atoms of
two scored residues
within 3.5 Å & the
perturbed
residues \\
Catalytic pKa (6) &
PROPKA pKa (min,
max, mean, range)
and pKa shift (mean,
max absolute) &
coordinates; PROPKA3
\cite{olsson2011,sondergaard2011}
pKa of the scored
ionisable residues;
the shift is the pKa
minus the pKa of the
model compound & the
perturbed
residues \\
Interface confidence
(2) & ipTM;
per-chain pTM
minimum & predictor
output of Chai-1
\cite{chai2024};
ipTM as introduced
for
AlphaFold-Multimer
\cite{evans2021};
the pTM of each
chain, from the
TM-score head of
AlphaFold
\cite{jumper2021} &
the protein--ligand
interface \\
Active-site accuracy
(2) & AME RMSD
(catalytic-atom RMSD
to a reference
structure after
backbone
superposition, best
of the predicted
models); ligand
clearance (closest
approach of the
ligand to the
backbone, best of
the models) &
predictor output and
a reference
structure; the
criteria of the
Atomic Motif Enzyme
benchmark of
RFdiffusion2
\cite{ahern2026}, as
implemented in this
project (equations
in 2.3) & the
catalytic atoms; the
ligand \\
Reference rows (6;
never counted as
metrics under test)
& resolution,
sequence length, net
charge, tyrosine
fraction, pLDDT, pTM
& deposit, sequence
or predictor; net
charge is (K + R) −
(D + E) counted in
the sequence; pLDDT
and pTM from
AlphaFold
\cite{jumper2021} &
whole protein \\
\end{longtable}

\textbf{Which
metrics are
reported.} The panel
contains 216 named
quantities, which
collapse to 191
distinct metrics
once literal aliases
are merged; 63 of
them are bookkeeping
counters and
constants on which
no test can be run.
Whether a metric can
be tested in an
experiment is
decided from what
the metric reads,
whether it is
defined under the
experimental
condition and the
design of the
experiment, never
from the result. The
main text reports
the 33 metrics that
are valid and tested
in all four tests
(Table 1), so that
one set of metrics
runs through the
whole paper, plus
six reference rows
chosen before any
result was seen. The
Chai-1 combined
score, a rescaled
ipTM and not an
independent metric,
is scored but kept
out of the main set.
The role of each of
the 191 distinct
metrics and the
reason for it are in
Table S1; what each
of the 216 named
quantities reads,
and over which
residues it is
scored, is in Table
S2; whether each can
be tested in each of
the 30 tests, and if
not why, is in
Tables S3 (one row
per metric and test)
and S4 (counts by
class of metric);
the counts of the
earlier whole-panel
analysis restated on
the eligible metrics
only are in Table
S5; and the list of
the tests, with
comparator, control
type and role, is in
Table S6.

\textbf{Notation.}
Every Results table
ranks a set of
\emph{items} (the
metrics, the
reference rows and,
in 2.4, the declared
baselines): \(m\) is
an item, \(M\) the
number of items
ranked in one table
and \(x_u(m)\) the
value of item \(m\)
on a unit \(u\) (a
structure, a
prediction, a
variant or a
design). Bootstrap
and permutation
tests use
\(N_{\mathrm{boot}} = N_{\mathrm{perm}} = 5{,}000\)
draws unless another
number is stated;
\(q_{\alpha}\) is
the
\(\alpha\)-quantile
of a set of values
and
\(\mathbf{1}(\cdot)\)
is 1 if its argument
is true and 0
otherwise. The
statistics (\(A_m\),
\(\rho_m\),
\(R_{m,s}\) and the
best-of-\(M\)
threshold from the
max-statistic
\(T^{(r)}\)) are
defined with their
equations in 2.2 to
2.5.

\begin{center}\rule{0.5\linewidth}{0.5pt}\end{center}

\subsection{2.2
Activity
discrimination test:
do the metrics
separate a dead
enzyme from an
active one?}

\textbf{Experiment.}
Every metric is
scored on real
proteins in which
one state is
catalytically dead
and the other
active, and the
metrics are ranked
by how well their
values separate the
two states. A
zymogen, for
example, is an
inactive precursor
that already
contains the
catalytic residues
of the mature
enzyme: a metric
that detected
catalytic competence
would separate the
two structures,
whereas a metric
that only reads the
geometry of the
catalytic residues
would not, because
that geometry is
almost the same in
both. Two datasets
are used, each with
its own table and
its own n: 21
zymogen--mature
pairs (2.2.1;
results in 3.1.1,
Table 5) and 192
plated designs, of
which 16 are active
(2.2.2; results in
3.1.2, Table 6). The
procedure is the
same for both:

\begin{enumerate}
\def\labelenumi{\arabic{enumi}.}
\tightlist
\item
  assemble the
  structures and
  label each one
  (dead or active;
  active or no
  active);
\item
  compute every item
  on every
  structure: the
  structure-space
  metrics on the
  deposited
  coordinates, the
  prediction-based
  metrics on a
  Chai-1
  \cite{chai2024}
  prediction of the
  structure together
  with its ligand;
\item
  compute the AUROC
  of every item and
  a bootstrap
  interval that
  resamples the
  clustering unit;
\item
  rank the items by
  the direction-free
  AUROC and compare
  the top of the
  ranking with the
  best-of-\(M\)
  permutation
  threshold.
\end{enumerate}

\subsubsection{2.2.1
Discriminate natural
active and dead
enzymes}

\textbf{Dataset.}
Pairs of structures
of the same protein,
the zymogen (dead)
and the mature
enzyme (active). The
structures were
downloaded from the
RCSB Protein Data
Bank
\cite{burley2025,berman2000};
the proteins and the
propeptide
boundaries come from
UniProt
\cite{uniprot2023}.
The pairs were
assembled in the
steps of Table 2,
and the evaluation
set is the last:

\begin{longtable}[]{@{}
  >{\raggedright\arraybackslash}p{(\columnwidth - 4\tabcolsep) * \real{0.2143}}
  >{\raggedright\arraybackslash}p{(\columnwidth - 4\tabcolsep) * \real{0.3214}}
  >{\raggedright\arraybackslash}p{(\columnwidth - 4\tabcolsep) * \real{0.4643}}@{}}
\caption*{\textbf{Table
2. Assembly of the
zymogen--mature
pairs.} From the
UniProt search to
the 21-pair
evaluation set of
Results 3.1.1; the
last row is the
evaluation
set.}\tabularnewline
\toprule\noalign{}
\begin{minipage}[b]{\linewidth}\raggedright
Step
\end{minipage} &
\begin{minipage}[b]{\linewidth}\raggedright
Cases
\end{minipage} &
\begin{minipage}[b]{\linewidth}\raggedright
Rule
\end{minipage} \\
\midrule\noalign{}
\endfirsthead
\toprule\noalign{}
\begin{minipage}[b]{\linewidth}\raggedright
Step
\end{minipage} &
\begin{minipage}[b]{\linewidth}\raggedright
Cases
\end{minipage} &
\begin{minipage}[b]{\linewidth}\raggedright
Rule
\end{minipage} \\
\midrule\noalign{}
\endhead
\bottomrule\noalign{}
\endlastfoot
UniProt discovery &
reviewed proteins of
EC 3.4.* with a
propeptide feature
and PDB structures &
a protein is
shortlisted when its
PDB chain ranges
give both a
candidate zymogen
and a candidate
mature X-ray
entry \\
Shortlisted proteins
& 100 & the proteins
that pass the rule
above \\
PDB entities
classified & 5,717
(1,182 zymogen,
4,528 mature, 7
indeterminate) &
each entity is
called zymogen or
mature from its
SIFTS
\cite{dana2019}
alignment to the
UniProt chain and
propeptide
boundaries, never
from the entry
title \\
Candidate pairs & 55
pairs of 55 proteins
& one zymogen and
one mature entity
per protein; both
X-ray, no engineered
mutation, resolution
at most 2.5 Å \\
Scored pairs (wider
set) & 49 &
structure-space
metrics scored on
both forms (S7);
trapping-mutant
pairs were examined
and retired \\
Re-predicted pairs &
28 & both forms
predicted with
Chai-1, three seeds
each (S8) \\
\textbf{Evaluation
set} & \textbf{21} &
the two forms share
a ligand; the 7 apo
pairs are
excluded \\
\end{longtable}

The 21 pairs are 21
proteins, all
peptidases (9
serine, 8 cysteine
and 4 aspartic; EC
3.4.21, 3.4.22,
3.4.23), with shared
ligands such as
inhibitors and metal
ions. In the wider
set of 49 pairs the
median
catalytic-constellation
RMSD between the two
forms is 0.207 Å
(0.484 Å after
backbone
superposition), so
the chemistry of the
site is intact in
the dead form. The
definitions of the
six pair sets used
in 3.1 (55, 49, 45,
28, 21 and 96 pairs)
and the table that
uses each are in
Table S9.

\textbf{Why this
dataset.} A zymogen
and its mature
enzyme hold the
catalytic residues
constant while the
enzyme is dead in
one form and active
in the other, so
they are a
\emph{real}
negative: the labels
are correct by
definition. They
answer the objection
that a metric may
only recognise
synthetic negatives.
The 21-pair subset
is the evaluation
set because it is
the only one on
which every main-set
metric,
structure-space and
prediction-based,
and every reference
row has a value on
both forms of every
pair, so that one n
applies to every row
of Table 5. The
pairs were selected
by ligand sharing
and not by any
metric value; the
choice was made
after the 49-pair
and 28-pair analyses
had been read, and
those analyses are
in the Supplement
(S7, S8). A second
class of dead
enzymes,
substrate-trapping
mutants, was
verified against the
primary literature
(26 of 30 were not
confirmed dead) and
retired (S10).

\textbf{Metrics
scored.} The 33
main-set metrics
(Table 1) and the 6
reference rows, 39
items in all, are
computed on both
forms of each pair.
The 29
structure-space
metrics are computed
on the deposited
coordinates over the
catalytic residues,
taken from M-CSA
\cite{ribeiro2018}
and UniProt
\cite{uniprot2023}
annotation. For the
4 prediction-based
metrics both forms
of each pair were
re-predicted with
Chai-1 together with
the shared ligand
(three seeds,
averaged), and
active-site accuracy
is measured against
each form's own
deposited structure.
Table S11 lists
every metric scored
on the 21 pairs
(158, including
whole-protein,
sequence-only and
PLACER
\cite{anishchenko2025}
comparators). The
reference rows are
reported in full,
with intervals, for
every activity test
in Table S12.

\textbf{Ranking
criteria.} Let \(P\)
be the dead forms
and \(Q\) the active
forms,
\(|P| = |Q| = n = 21\).
The AUROC of item
\(m\) is the
probability that a
randomly chosen dead
form has a larger
value than a
randomly chosen
active form, ties
counting one half
\cite{hanley1982}:

\begin{equation*}A_m \;=\; \frac{1}{|P|\,|Q|}\sum_{p \in P}\sum_{q \in Q}\Big[\,\mathbf{1}\big(x_p(m) > x_q(m)\big) + \tfrac{1}{2}\,\mathbf{1}\big(x_p(m) = x_q(m)\big)\Big] \tag{1}\end{equation*}

This equals the
Mann--Whitney
statistic \(U\)
divided by
\(|P|\,|Q|\)
\cite{mann1947}.
\(A_m = 0.5\) means
that the item does
not separate the two
forms, \(A_m = 1\)
that every dead form
is above every
active form and
\(A_m = 0\) that
every dead form is
below every active
form. We ask whether
an item separates
the forms, not in
which direction, so
items are ranked by

\begin{equation*}A^{*}_m \;=\; \max\!\big(A_m,\; 1 - A_m\big) \;\ge\; 0.5 , \tag{2}\end{equation*}

and the directed
\(A_m\) is kept so
that the direction
is visible (``higher
in'' in Table 5 is
the form in which
the item takes
larger values: dead
if \(A_m > 0.5\)).
Because \(A^{*}_m\)
cannot fall below
0.5, the lower end
of its interval is
close to 0.5 for
almost every item
and is not
informative by
itself.

The interval is a
percentile bootstrap
that resamples
\textbf{pairs} with
replacement, so that
both forms of a
drawn pair enter
together
\cite{efron1979,field2007}.
With
\(b = 1,\dots,N_{\mathrm{boot}}\)
draws (5,000, seed
20260806),
\(A^{*(b)}_m\) is
recomputed on each
resampled set and

\begin{equation*}\Big[\,q_{0.025}\big(A^{*(1)}_m,\dots,A^{*(N_{\mathrm{boot}})}_m\big),\;\; q_{0.975}\big(A^{*(1)}_m,\dots,A^{*(N_{\mathrm{boot}})}_m\big)\Big] \tag{3}\end{equation*}

is the 95\%
interval. A pair is
also a protein, so
the design-level and
the cluster-level N
are both 21.

Table 5 ranks
\(M = 39\) items,
and the top of such
a ranking looks good
by chance. To say
how good, the labels
are permuted. Within
a pair the two forms
are exchangeable if
the item is
unrelated to the
label, so in
replicate
\(r = 1,\dots,N_{\mathrm{perm}}\)
we draw
\(s_i^{(r)} \sim \mathrm{Bernoulli}(1/2)\)
for every pair
\(i\), swap the dead
and active labels of
the pair when
\(s_i^{(r)} = 1\),
and recompute
\(A^{*(r)}_m\) for
every item. The same
swaps serve all
items, which keeps
the correlation
between metrics. The
largest value over
the items in each
replicate,

\begin{equation*}T^{(r)} \;=\; \max_{m = 1,\dots,M} A^{*(r)}_{m}, \tag{4}\end{equation*}

has the distribution
of the best AUROC of
\(M\) unrelated
items. Its median
and its 95th
percentile
\(q_{0.95}(T)\) are
the
\textbf{best-of-39
threshold} (median
0.63, 95th
percentile 0.68): an
observed \(A^{*}_m\)
above it is higher
than the best of 39
unrelated items in
95\% of simulations.
Taking the maximum
over items is the
max-statistic form
of permutation
control of multiple
comparisons
\cite{westfall1993}.
The threshold is a
reference level for
reading the ranking.
For the complete
table of all 158
metrics (S11) the
same procedure gives
the best-of-158
threshold. In Table
5 the AUROCs above
the threshold are in
bold. No threshold
is adopted for this
set: we report
whether any main-set
metric exceeds the
best-of-39
threshold, and no
point estimate is
offered as a
headline when its
interval is wider
than 0.2.

\subsubsection{2.2.2
Discriminate de novo
designed active and
no active enzymes}

\textbf{Dataset.} A
deposited
metallohydrolase
design campaign
\cite{kim2026}: 192
designs on two
96-well plates (two
campaigns of 96
designs, the two
96-design zinc
campaigns that are
also reported with
RFdiffusion2
\cite{ahern2026}).
Of the 192,
\textbf{16 have a
reported kcat/KM}
(\emph{active}) and
\textbf{176 are no
active}. The designs
lie in \textbf{136
backbone clusters},
and the 16 active
designs lie in 8 of
them (5 in each of
two clusters and one
in each of six
others). The
catalytic motif is
the three histidines
that coordinate the
zinc, derived
geometrically and
checked against 23
curated annotations
(23 of 23 agree).
For the predictions
the ligand is the
design's
transition-state
analogue together
with the zinc. One
set serves every row
of Table 6: every
item is scored on
all 192 designs.

\textbf{Why this
dataset.} These are
the deposited de
novo enzyme designs
for which models and
measured activity
are both available,
so they let the
metrics be tested on
designs and not only
on natural enzymes.
A design absent from
the table of
characterised hits
was not reported as
active, and the
comparison is active
against no active.
The plate analyses
were planned as
descriptive, without
an AUROC, before the
data were analysed;
the descriptive
analyses (rank
correlation among
the 16, hits per
plate) are reported
in 2.4.2 and the
AUROC, added later,
in this subsection.

\textbf{Metrics
scored.} The 36
items are the 29
structure-space
main-set metrics,
the two
prediction-based
main-set metrics
that the plates
define (ipTM and the
per-chain pTM
minimum) and five
reference rows
(tyrosine fraction,
net charge, sequence
length, pLDDT and
pTM; designs have no
crystallographic
resolution). An item
must have a value on
at least 80\% of the
designs. The AME
RMSD and the ligand
clearance are
measured on the
plates against each
design's own model,
which is a different
metric from AME
against a deposited
structure (the two
flavours are never
pooled), so they are
in the Supplement
and not among the
ranked items. Table
S13 lists every
metric scored on the
designs (128).

\textbf{Ranking
criteria.} The
statistic is Eq. (1)
and (2) with \(P\)
the 16 active
designs and \(Q\)
the 176 no active
designs, over the
192 designs;
``higher in'' in
Table 6 is the group
in which the item
takes larger values.
The interval
resamples the
\(G = 136\)
\textbf{backbone
clusters} with
replacement, all
designs of a drawn
cluster entering
together, so that
designs built on one
backbone are not
counted as
independent (5,000
draws, seed
20260807; a draw
with no active
design has no AUROC
and is dropped):

\begin{equation*}\Big[\,q_{0.025}\big(A^{*(b)}_m\big),\;\; q_{0.975}\big(A^{*(b)}_m\big)\Big], \qquad A^{*(b)}_m = \max\!\big(A^{(b)}_m,\, 1 - A^{(b)}_m\big), \tag{5}\end{equation*}

with \(A^{(b)}_m\)
from Eq. (1) on the
designs of the drawn
clusters. The
design-level N is
192 and the
cluster-level N is
136, with 16 actives
in 8 clusters. The
chance reference is
the permutation of
Eq. (4) with the
activity labels
permuted among the
192 designs instead
of swapped within
pairs, \(M = 36\)
items and 5,000
permutations: the
\textbf{best-of-36
threshold} (median
0.66, 95th
percentile 0.73),
and the best-of-128
threshold (0.75) for
the 128 metrics of
Table S13. Because
10 of the 16 active
designs lie in two
backbone clusters,
designs are not
exchangeable and the
threshold
understates the
chance range. Table
6 marks the AUROCs
above the best-of-36
threshold in bold.

\begin{center}\rule{0.5\linewidth}{0.5pt}\end{center}

\subsection{2.3
Substrate
discrimination test:
do prediction-based
metrics know which
ligand they were
given?}

\textbf{Experiment.}
A structure
predictor (Chai-1
\cite{chai2024}) is
given a protein
together with its
own substrate, and
then together with a
wrong one, and the
prediction-based
quantities are
ranked by how well
they separate the
cognate prediction
from a wrong-ligand
prediction of the
same enzyme. For
example, Chai-1
predicts a hydrolase
together with its
own substrate and
again together with
a substrate of a
different enzyme
class; if ipTM falls
when the wrong
substrate is given,
the metric ``knows''
which substrate
belongs to the
protein, and the
AUROC measures how
reliably it does.
Chai-1 is run on the
sequence, the ligand
(as a chemical
structure) and,
where relevant, the
catalytic metal,
with default
settings and a fixed
seed; each
prediction yields
five models and
their confidence
scores. Only the
ligand supplied
changes; the protein
sequence is the same
in every condition,
and four conditions
are built for every
enzyme (results in
3.2, Table 7):

\begin{enumerate}
\def\labelenumi{\arabic{enumi}.}
\tightlist
\item
  \textbf{cognate}:
  the substrate of
  the enzyme (three
  seeds, 101, 202
  and 303; the
  cognate condition
  has two further
  seeds that serve
  as the noise floor
  and are not used
  in the
  evaluation);
\item
  \textbf{same
  class}: a
  different
  substrate of an
  enzyme of the same
  enzyme class;
\item
  \textbf{different
  class}: a
  substrate of an
  enzyme of a
  different class;
\item
  \textbf{decoy}: a
  ligand taken from
  an unrelated
  structure.
\end{enumerate}

Wrong ligands have
three seeds each,
and the catalytic
metal is kept in
every condition. The
apo condition (no
ligand) is excluded
for interface and
ligand-dependent
terms, because
without a ligand
there is no
interface to be
confident about.

\textbf{Dataset.}
Natural enzymes of
M-CSA
\cite{ribeiro2018}
with a bound cognate
ligand: 59 enzymes
enter the
experiment, and the
\textbf{evaluation
set is the 55
natural enzymes} on
which every scored
quantity has a
defined value in the
cognate condition
and in all three
wrong-ligand
conditions, and on
which every ligand
that was supplied
was present in the
prediction. The four
exclusions are
enzymes whose
cognate or wrong
ligand contains a
heme-type group that
cannot be built into
a chemical
structure, so the
predictor drops it
and the prediction
is effectively apo
for that cell. The
rule is about the
availability of
values and not about
any result. The
experiment is
balanced: seeds 101,
202 and 303 in every
condition.
Twenty-eight of the
59 enzymes are also
in the 143-enzyme
set of 2.5.

\textbf{Why this
dataset.} A wrong
ligand can only be
defined for an
enzyme whose
substrate is known,
and the M-CSA
entries with the
cognate reaction
ligand present (150
of 434) are the
natural enzymes for
which this holds.
Natural enzymes
carry the main
analysis because
their active-site
accuracy is measured
against the
deposited structure;
the 30 de novo
designs, whose
active-site accuracy
is measured against
the design's own
model and is never
pooled with it, are
in S14. The choice
of the 55 systems
was made after the
earlier all-system
analysis had been
read.

\textbf{Metrics
scored.} Every
prediction is scored
on 33
prediction-based
quantities, of which
Table 7 ranks 30.
The three summaries
of the Chai-1
combined score are
scored but not
listed, because the
combined score is
not a main-set
metric (see ``Which
metrics are
reported''); their
AUROCs are in S14.
The 30 listed
quantities are the
four main-set
metrics of Table 1,
13 whole-prediction
confidence summaries
(pLDDT, pTM, the
per-chain pTM of the
protein chain and
the inter-chain
clash flag, two of
which, the mean
pLDDT and mean pTM,
are the reference
rows), the AME pass
bit, and four
bookkeeping counters
(models, aligned and
catalytic atoms,
symmetry swaps) that
serve as negative
controls. Each score
\(s_k\) is returned
for each of the
\(K = 5\) models,
and every score is
summarised over the
models as

\begin{equation*}\bar s = \frac{1}{K}\sum_{k=1}^{K} s_k, \qquad s_{\max} = \max_{k} s_k, \qquad s_{\mathrm{SD}} = \sqrt{\frac{1}{K-1}\sum_{k=1}^{K}\big(s_k - \bar s\big)^{2}}, \tag{6}\end{equation*}

and the RMSD also as
the minimum over the
models, so that one
metric appears in
several rows of
Table 7, one per
summary. Two
quantities need a
definition of their
own (the other
summaries of
whole-prediction
confidence are those
of Chai-1
\cite{chai2024}).

\textbf{AME RMSD}
(active-site
accuracy; the
criterion of the
Atomic Motif Enzyme
benchmark
\cite{ahern2026}). A
model is superposed
on the reference
structure by the
Kabsch algorithm on
the C\(\alpha\)
atoms of all
residues shared with
the reference (the
whole shared
backbone, and not
the catalytic
residues alone,
which would let the
site translate
freely inside the
fold and deflate the
RMSD). The RMSD of
the catalytic atoms
is then

\begin{equation*}\mathrm{RMSD}_k \;=\; \sqrt{\frac{1}{L}\sum_{l=1}^{L}\big\lVert \mathbf{y}_l - \mathbf{x}_{l}^{(k)}\big\rVert^{2}}, \qquad \mathrm{AME\ RMSD} = \min_{k} \mathrm{RMSD}_k, \tag{7}\end{equation*}

where
\(l = 1,\dots,L\)
runs over the heavy
atoms of the
catalytic residues
that keep their
residue type in the
prediction,
\(\mathbf{y}_l\) is
the atom in the
reference,
\(\mathbf{x}^{(k)}_l\)
the superposed atom
of model \(k\), and
symmetry-equivalent
atoms (the two
oxygens of a
carboxylate, for
example) are
assigned to give the
smaller RMSD. The
reference is the
deposited structure
for natural enzymes
and the design's own
model for designs;
the two flavours are
never pooled. The
variant without
symmetry resolution
is also listed (``no
symmetry
resolution'' in
Table 7).

\textbf{Ligand
clearance} is the
closest approach of
the ligand to the
protein backbone
atoms N, C\(\alpha\)
and C, best case
over the models:

\begin{equation*}C_k \;=\; \min_{a \in \text{ligand}}\ \min_{b \in \{\mathrm{N},\,\mathrm{C}\alpha,\,\mathrm{C}\}} \big\lVert \mathbf{x}_a^{(k)} - \mathbf{x}_b^{(k)}\big\rVert, \qquad \mathrm{clearance} = \max_{k} C_k , \tag{8}\end{equation*}

so that a larger
value means more
room for the ligand
(the AME criterion
is \(C > 1.5\) Å).
Interface confidence
(ipTM, per-chain pTM
minimum) and ligand
clearance need a
ligand and are
undefined without
one.

\textbf{Ranking
criteria.} For each
quantity \(m\) and
each kind of wrong
ligand
\(c \in \{\text{same class},\ \text{different class},\ \text{decoy}\}\),
the AUROC of Eq. (1)
is computed with the
cognate predictions
as the positive
class and the
wrong-ligand
predictions of kind
\(c\) as the
negative class, one
value per system
(the mean over its
three seeds), and
made direction-free
as in Eq. (2);
quantities are
ranked by the mean
of the three,

\begin{equation*}\bar A_m \;=\; \frac{1}{3}\sum_{c}\, A^{*}_{m,c}. \tag{9}\end{equation*}

Each \(A^{*}_{m,c}\)
has the percentile
bootstrap interval
of Eq. (3) that
resamples systems
(5,000 draws, seed
20260806); ``higher
in'' is the
condition in which
the quantity takes
larger values. The
chance reference is
the max-statistic
permutation of Eq.
(4): in each of
5,000 replicates the
four condition
labels (cognate,
same class,
different class,
decoy) are permuted
at random within
every system, the
three AUROCs and
their mean are
recomputed for every
quantity, and
\(T^{(r)} = \max_{m} \bar A^{(r)}_m\)
over the \(M = 30\)
ranked quantities;
its median and 95th
percentile are the
\textbf{best-of-30
threshold} (median
0.56, 95th
percentile 0.61)
\cite{westfall1993}.
Table 7 marks in
bold the mean AUROCs
above it. No
threshold is
adopted; we report
which quantities
exceed it.

To check whether a
ligand-size effect
explains the
discrimination, the
cells (a cognate
ligand and one wrong
ligand of an enzyme)
are stratified by
the difference in
heavy-atom count,
\(\Delta h = h_{\text{cognate}} - h_{\text{wrong}}\),
and in charge, and
the paired win
fraction of the
quantity is reported
(S15):

\begin{equation*}W \;=\; \frac{1}{n_{\mathrm{cells}}}\sum_{\mathrm{cells}}\Big[\mathbf{1}\big(x_{\mathrm{cognate}} > x_{\mathrm{wrong}}\big) + \tfrac12\,\mathbf{1}\big(x_{\mathrm{cognate}} = x_{\mathrm{wrong}}\big)\Big]. \tag{10}\end{equation*}

\begin{center}\rule{0.5\linewidth}{0.5pt}\end{center}

\subsection{2.4
Activity ranking
test: do the metrics
track measured
activity?}

\textbf{Experiment.}
We ask whether the
order of real
variants (or real
designs) by a metric
agrees with their
order by measured
activity, and
whether the metric
does better than
simply knowing how
close the change is
to the active site.
For example, the
variants are lined
up by how strongly
the assay found each
one impaired and
again by how much a
metric changes for
each one; if the two
orders agree, the
metric ranks the
variants as the
assay does, and the
rank correlation
measures how well.
The baselines ask
whether knowing only
the distance from
the mutated residue
to the catalytic
glutamates ranks
them as well. This
is a ranking claim,
so the statistic is
a rank correlation;
the AUROC of the
binary contrast
(impaired against
wild-type-like
variants) is given
for BglB in the
Supplement (S16).
Two datasets are
used, each with its
own table and its
own n: 432 BglB
variants (2.4.1;
results in 3.3.1,
Table 8) and the 16
plated designs that
have a measured
activity (2.4.2;
results in 3.3.2,
Table 9).

\subsubsection{2.4.1
Ranking natural
enzyme activity}

\textbf{Dataset.}
β-glucosidase B
(BglB) from
\emph{Paenibacillus
polymyxa},
characterised by the
D2DCure consortium
\cite{carlin2016}
(data retrieved from
d2dcure.com on
2026-08-06; the
licence is unstated,
so variant-level
data are not
redistributed).
Variant numbering is
offset by three from
the numbering of the
crystal structure
(RCSB PDB entry 2JIE
\cite{burley2025,berman2000};
structure position =
variant position +
3). The catalytic
residues are Glu167
(proton donor) and
Glu356 (nucleophile)
in structure
numbering, taken
from the UniProt
\cite{uniprot2023}
active-site
features, the
conserved motifs
TINEP and ITENG, and
the covalent link to
the glucosyl
intermediate in the
deposited structure.
The funnel from the
measured rows to the
evaluation set is in
Table 3:

\begin{longtable}[]{@{}
  >{\raggedright\arraybackslash}p{(\columnwidth - 4\tabcolsep) * \real{0.2075}}
  >{\raggedright\arraybackslash}p{(\columnwidth - 4\tabcolsep) * \real{0.3396}}
  >{\raggedright\arraybackslash}p{(\columnwidth - 4\tabcolsep) * \real{0.4528}}@{}}
\caption*{\textbf{Table
3. From the D2DCure
rows to the
432-variant
evaluation set.} The
funnel for Results
3.3.1; the last row
is the set on which
every row of Table 8
is
computed.}\tabularnewline
\toprule\noalign{}
\begin{minipage}[b]{\linewidth}\raggedright
Step
\end{minipage} &
\begin{minipage}[b]{\linewidth}\raggedright
Cases
\end{minipage} &
\begin{minipage}[b]{\linewidth}\raggedright
Rule
\end{minipage} \\
\midrule\noalign{}
\endfirsthead
\toprule\noalign{}
\begin{minipage}[b]{\linewidth}\raggedright
Step
\end{minipage} &
\begin{minipage}[b]{\linewidth}\raggedright
Cases
\end{minipage} &
\begin{minipage}[b]{\linewidth}\raggedright
Rule
\end{minipage} \\
\midrule\noalign{}
\endhead
\bottomrule\noalign{}
\endlastfoot
Measured rows &
1,288 rows, of which
272 are wild-type
replicates measured
at eight
institutions & the
D2DCure file
(latin-1 encoded) \\
Single-point
variants & 655 at
207 positions & the
variants of the
characterisation
data \\
Set aside & 108
destabilised, 83
non-expressed, 32
with insufficient
data & a variant is
\textbf{folded} if
it was expressed and
is not destabilised
(T50 or Tm less than
2 standard
deviations below
wild type); the
kinetics of the
others cannot be
attributed to the
catalytic site \\
Folded variants with
kinetics & 432 at
175 positions: 166
impaired, 254
wild-type-like, 12
improved &
\textbf{impaired} if
kcat/KM is at least
2 standard
deviations (a
3.8-fold loss or
more) below wild
type,
\textbf{wild-type-like}
if within 2 standard
deviations,
\textbf{improved}
above \\
\textbf{Evaluation
set} & \textbf{432
variants} & one set
serves every row of
Table 8; a
structure-space
metric has a value
for 431 of the 432
variants and a
prediction-based
metric for all 432;
the dataset's own
Rosetta score exists
for 248 \\
\end{longtable}

The spread of the
wild-type replicates
(0.293 log10) is the
noise floor, and
each variant is
compared with the
wild type of the
same institution.
The binary contrast
of the Supplement
(S16) uses the 420
variants that are
impaired or
wild-type-like.
Substrate: the
prediction-based
metrics use the
assay substrate
4-nitrophenyl
β-D-glucopyranoside
(pNPG), the reporter
substrate of the
published BglB
assays
\cite{carlin2016,carlin2017},
and no metal,
because BglB is
metal-independent.

\textbf{Why this
dataset.} BglB is
the powered
calibration set of
natural-enzyme
variants: hundreds
of single-point
variants of one
enzyme with
kinetics, expression
and thermal
stability (T50, Tm)
measured in one
consortium, so that
a ranking claim can
be tested at a scale
that no set of
designs offers, and
so that a loss of
activity from a
failure to fold (a
destabilised or
non-expressed
variant) can be
separated from a
loss at the
catalytic site. The
analysis plan was
committed and dated
before any mutant
structure was built
and before any
metric was joined to
a label.

\textbf{Metrics
scored.} Each
variant is built in
silico with the same
rotamer-space
substitution, repack
and minimisation as
in 2.5 (the wild
type goes through
the identical
treatment with five
seeds), and the
metrics are computed
over the two
catalytic glutamates
for every variant.
The prediction-based
metrics are computed
on Chai-1
predictions of each
variant and of the
wild type with the
assay substrate; AME
is measured against
the wild-type
crystal, which
carries a covalent
glucosyl
intermediate and not
the substrate. The
41 ranked items of
Table 8 are the 28
structure-space
main-set metrics
that BglB defines
(the omega term does
not change under a
side-chain
substitution on a
fixed backbone), the
4 prediction-based
main-set metrics,
four reference rows
(tyrosine fraction,
net charge, pLDDT
and pTM; the
sequence length is
constant over single
substitutions and
cannot be ranked)
and five
\textbf{declared
baselines}, fixed
before any metric
was joined to a
label, each pointed
so that a larger
value means more
impaired: the
absolute change in
side-chain volume;
the negative
distance from the
mutated residue to
the nearest
catalytic residue in
the wild-type
structure (the
\emph{distance
baseline}); the
negative relative
solvent-accessible
area of the mutated
residue (burial);
the negative
BLOSUM62 score of
the substitution
\cite{henikoff1992};
and the Rosetta
score that comes
with the dataset
(available for 248
of the variants, so
compared on that
subset).

\textbf{Ranking
criteria.} The
measured
\textbf{impairment}
of variant \(v\) is
the log-ratio of its
kcat/KM to the
wild-type value of
the same
institution,

\begin{equation*}I_v \;=\; \log_{10}\frac{\kappa_{\mathrm{wt}}}{\kappa_{v}}, \tag{11}\end{equation*}

larger for a more
impaired variant
(\(\kappa\) is
kcat/KM). A metric
enters as the size
of its change, the
absolute difference
between the variant
and the treated
unsubstituted
protein,
\(\lvert\Delta_v(m)\rvert = \lvert m(v) - m(\mathrm{wt})\rvert\),
and a declared
baseline enters as
its own score. The
direction is fixed
in advance (a larger
change in a more
impaired variant),
so the directed
Spearman correlation
is reported and a
negative value is
kept as it is. With
\(r_v\) the rank of
the item's score and
\(s_v\) the rank of
\(I_v\) among the
\(n = 432\) variants
(ties take the
average rank), the
Spearman rank
correlation is the
Pearson correlation
of the ranks
\cite{spearman1904}:

\begin{equation*}\rho_m \;=\; \frac{\sum_{v}(r_v-\bar r)(s_v-\bar s)}{\sqrt{\sum_{v}(r_v-\bar r)^2\;\sum_{v}(s_v-\bar s)^2}} \tag{12}\end{equation*}

\(\rho = 1\) means
that the item orders
the variants exactly
as the assay does,
\(\rho = 0\) that
the two orders are
unrelated, and
\(\rho = -1\) that
they are opposite.
The 95\% interval is
a percentile
bootstrap (Eq. 3
with \(\rho\) in
place of \(A^{*}\))
that resamples the
175
\textbf{positions}
with replacement
(5,000 draws, seed
20260807), because
variants at one
position are not
independent; the
design-level N is
432 and the
cluster-level N is
175. To say whether
an item is better
than the distance
baseline, the
difference is
computed on the same
draws,

\begin{equation*}D_m^{(b)} \;=\; \rho_m^{(b)} - \rho_{\mathrm{dist}}^{(b)}, \tag{13}\end{equation*}

and the item
\textbf{beats} the
baseline only if the
95\% interval of
\(D_m\),
\(q_{0.025}(D_m^{(b)})\)
to
\(q_{0.975}(D_m^{(b)})\),
lies above 0. Table
8 ranks \(M = 41\)
items by \(\rho_m\),
and the top of such
a ranking looks good
by chance, so the
measured impairments
are permuted among
whole positions that
have the same number
of variants (a block
permutation, so that
the variants at one
position keep their
impairments
together, as in the
bootstrap),
\(\rho^{(r)}_m\) is
recomputed for every
item with the same
permutation, and

\begin{equation*}T^{(r)} \;=\; \max_{m = 1,\dots,M} \rho^{(r)}_{m}, \qquad r = 1,\dots,N_{\mathrm{perm}} . \tag{14}\end{equation*}

The median and the
95th percentile of
\(T\) are the
\textbf{best-of-41
threshold}: the
correlation that the
best of 41 unrelated
items would exceed
in 5\% of
permutations (median
0.24, 95th
percentile 0.32; for
the complete table
S16, over all 114
analysed quantities,
0.32). The same
procedure applies to
the declared
baselines, whose raw
score takes the
place of
\(\lvert\Delta\rvert\).
It is a reference
level for reading
the ranking; a
\(\rho\) above it is
in bold in Table 8.
The signed
correlation of the
pre-specified plan
(S17) and the AUROC
of the binary
contrast (S16) are
given beside it. The
ten sensitivity
analyses (variants
near to and distal
from the catalytic
residues, restricted
subsets of variants
or metrics, other
thresholds for the
labels, and each
institution left out
in turn) are
summarised in Table
S18, the 28
prediction-based
metrics, with their
own correlation and
AUROC columns, are
in Table S19, and
the declared
baselines are
reported in full in
Table S12.

\subsubsection{2.4.2
Ranking de novo
designed enzyme
activity}

\textbf{Dataset.}
The 192-design
campaign of 2.2.2
\cite{kim2026}. Only
the 16 designs with
a reported kcat/KM
have a measured
activity that an
item can be
correlated with, so
the
\textbf{evaluation
set is those 16
designs} and one n
serves every row of
Table 9. They lie in
8 of the 136
backbone clusters,
so the effective
number of
independent designs
is smaller than 16.
The 192-design plate
is used for the
AUROC of 2.2.2
(Table 6) and for
the filter-style
analysis (S20).

\textbf{Why this
dataset.} The same
designs as 2.2.2, so
that the
discrimination
(active against the
rest of the plate)
and the ranking
(among the actives)
are two readings of
one campaign. The
question a designer
asks of a metric is
whether it would
have enriched the
hits on a plate, and
this is a deposited
design campaign on
which that can be
asked. The plate
analyses of this
subsection are
\textbf{descriptive}
and were planned
before these data
were analysed; the
AUROC, added later,
is reported in
2.2.2.

\textbf{Metrics
scored.} The 36
items of Table 9 are
those of Table 6:
the 29
structure-space
main-set metrics,
the two
prediction-based
main-set metrics
that the plates
define (ipTM and the
per-chain pTM
minimum) and five
reference rows
(tyrosine fraction,
net charge, sequence
length, pLDDT and
pTM). Designs are
scored as deposited,
and their
predictions use the
design's
transition-state-analogue
ligand and zinc,
with accuracy
measured against the
design's own model.
AME and the ligand
clearance are in the
Supplement (S21) and
not among the ranked
items, as in 2.2.2.

\textbf{Ranking
criteria.} The
correlation is Eq.
(12) between the
native value of the
item on the design
and
\(\log_{10}\mathrm{kcat/KM}\)
of the \(n = 16\)
designs that have a
reported value; no
direction is fixed
in advance, so items
are ranked by
\(\lvert\rho_m\rvert\)
and the sign is
reported (``higher
value means higher
or lower
activity''). The
interval is a 90\%
percentile bootstrap
that resamples
designs (2,000
draws). Beside it,
the correlation is
also computed with
the sequence length
partialled out of
both sides (S22,
S21): with the ranks
\(r_x\), \(r_y\) and
\(r_z\) of the item,
of the activity and
of the length,
centred, the
residual of each
rank on \(r_z\) is
taken and the
Pearson correlation
of the residuals is
the partial
correlation,

\begin{equation*}e_x = r_x - \frac{r_z^{\top} r_x}{r_z^{\top} r_z}\, r_z, \qquad e_y = r_y - \frac{r_z^{\top} r_y}{r_z^{\top} r_z}\, r_z, \qquad \rho_{xy\cdot z} = \mathrm{corr}(e_x, e_y). \tag{15}\end{equation*}

A \(\rho\) is in
bold in Table 9 when
its 90\% interval
excludes 0, and
about 10\% of
unrelated items are
expected to show
this by chance (0.1
× 36 = 3.6 of the
36). The
filter-style
analysis fills a
96-well plate only
from the quarter of
the 192 designs that
the item ranks
highest, or lowest:
with
\(n_{\mathrm{keep}} = 48\)
designs kept, of
which \(k\) are
active, the expected
hits per plate are

\begin{equation*}\text{hits per plate} \;=\; 96\,\frac{k}{n_{\mathrm{keep}}} \;=\; 2k , \tag{16}\end{equation*}

against
\(96 \times 16/192 = 8.0\)
for an unfiltered
plate. The band of a
random filter is the
distribution of
\(k\) when 48
designs are chosen
at random,
hypergeometric with
16 active designs
among 192,

\begin{equation*}\Pr(k \ge h) \;=\; \sum_{j \ge h}\frac{\binom{16}{j}\binom{176}{48-j}}{\binom{192}{48}} , \tag{17}\end{equation*}

so that 95\% of
random choices hold
at most 14.0 hits
per plate (\(h = 7\)
actives), a random
choice holds more
with probability
0.022 and either end
of a random item
does with
probability 0.044
(S20). The
correlation of every
structure-space
metric and
comparator, with and
without the sequence
length partialled
out, is in Table
S22, and the
enrichment at keep
fractions of 5, 10,
25 and 50\% in both
directions is in
Table S23
(structure-space
metrics) and Table
S24
(prediction-based
metrics).

\begin{center}\rule{0.5\linewidth}{0.5pt}\end{center}

\subsection{2.5
Detection ability
test: do the metrics
respond more to a
catalytic lesion
than to a matched
control?}

\textbf{Experiment.}
We damage the
catalytic residues
of natural enzymes,
in the structure for
the structure-space
metrics and in the
sequence of a
re-predicted complex
for the
prediction-based
metrics, and ask
whether each metric
moves more than it
does when the same
damage is made at a
matched
non-catalytic site
of the same protein.
We call each such
deliberate change a
\emph{lesion}. In a
serine protease, for
example, the
catalytic triad is
Asp, His and Ser.
The \emph{Ala step}
builds a structure
in which those three
residues are
replaced by Ala. The
\emph{control}
builds a structure
in which three
buried,
non-catalytic
residues of the same
original types are
replaced by Ala,
with the same
repacking. A metric
is \emph{specific}
if it changes
clearly more in the
first structure than
in the second. Two
kinds of lesion are
applied, each with
its own datasets and
table: the catalytic
lesion (the
chemistry ladder,
2.5.1; results in
3.4.1, Table 10) and
the second-shell
lesion (2.5.2;
results in 3.4.2,
Table 11).

\textbf{Change.}
Every change is made
in side-chain
rotamer space on the
fixed backbone,
followed by an 8 Å
fixed-backbone
repack with the
changed residues
held and constrained
minimisation
(PyRosetta
\cite{chaudhury2010});
protonation is
fixed. The
\textbf{chemistry
ladder} replaces
each catalytic
residue in one of
four steps of
increasing change in
side-chain volume:
\textbf{isosteric}
(for example
Ser→Cys, Asp→Asn,
Thr→Val; median
change 15.3 ų),
\textbf{non-isosteric}
(36.4 ų),
\textbf{to Ala}
(53.2 ų) and
\textbf{to Gly}
(80.5 ų).
\textbf{Second-shell
removal} mutates 1,
2 or 4 residues in
the shell around the
catalytic residues.
The
\textbf{detection
floor} replaces all
catalytic residues
by Gly at once,
which is the same
lesion as the Gly
step, and asks
whether the metric
notices. Table S25
gives the result of
the detection floor
and of the other
sanity-floor
conditions by class
of metric, and
checks that the
detection floor and
the Gly step are
identical. For the
prediction-based
metrics the same
four steps of the
chemistry ladder are
applied in the
\textbf{sequence}:
one catalytic
residue is
substituted per
prediction and the
complex is
re-predicted with
Chai-1
\cite{chai2024} (one
seed, 101).

\textbf{Comparator.}
Two comparisons are
made, and they
answer different
questions.

\begin{itemize}
\tightlist
\item
  \emph{Against the
  structure's own
  reference.} Every
  perturbed
  structure and its
  reference go
  through the
  identical
  treatment, because
  26 of 103 metrics
  move under the
  treatment alone;
  each arm is
  compared with its
  own treated
  reference.
\item
  \emph{Against the
  matched control.}
  The control makes
  the
  \textbf{identical
  substitution}
  (same original and
  same new residue
  type) at buried,
  non-catalytic
  positions of the
  same protein,
  matched on burial
  (within 0.10
  relative
  solvent-accessible
  area) and, in the
  primary analysis,
  on local packing,
  secondary
  structure and
  distance to the
  ligand where
  available. When no
  exact match
  exists, the
  matcher relaxes
  its criteria in a
  fixed order and
  records the
  loosest tier used
  (the loosest tier
  needed for each
  lesion test is in
  Table S26). Only
  comparisons where
  both arms change
  the same number of
  residues are
  counted (the
  balance of the
  numbers of
  residues changed
  in the two arms is
  in Table S27). For
  the re-predicted
  ladder each
  substitution is
  paired with the
  identical
  substitution at a
  buried
  non-catalytic
  position (same
  original and same
  new residue,
  matched on
  burial), and the
  change is measured
  against the
  unperturbed
  prediction of the
  same system at the
  same seed.
\end{itemize}

\textbf{Response.}
For an item \(m\),
an enzyme \(i\) and
an arm
\(a \in \{\mathrm{cat}, \mathrm{ctl}\}\)
(the catalytic arm
and the control
arm), the response
at step \(s\) is the
change of the item
from its own treated
reference,

\begin{equation*}\Delta^{a}_{i}(m, s) \;=\; m\big(\text{changed}^{a}_{i,s}\big) - m\big(\text{reference}^{a}_{i}\big) , \tag{18}\end{equation*}

summarised as the
median over enzymes.
For the
structure-space
metrics the
site-scoped items
are computed over
the changed residues
in each arm (the
catalytic residues
in the catalytic
arm, the control
residues in the
control arm); for
second-shell removal
both arms are scored
over the catalytic
residues. For the
prediction-based
metrics, both pair
members are scored
over the catalytic
set \emph{minus} the
substituted
position, so the
active-site accuracy
reports only the
displacement of the
remaining catalytic
residues. An item
\textbf{responds} at
a step when the 95\%
interval of its
median response
excludes zero.

\subsubsection{2.5.1
Detecting catalytic
lesion}

\textbf{Dataset.}
The datasets of this
test are listed in
Table 4:

\begin{longtable}[]{@{}
  >{\raggedright\arraybackslash}p{(\columnwidth - 4\tabcolsep) * \real{0.2407}}
  >{\raggedright\arraybackslash}p{(\columnwidth - 4\tabcolsep) * \real{0.3333}}
  >{\raggedright\arraybackslash}p{(\columnwidth - 4\tabcolsep) * \real{0.4259}}@{}}
\caption*{\textbf{Table
4. Datasets of the
catalytic lesion
test.} The sets on
which the catalytic
lesion is scored in
Results 3.4.1, the
sets added to them
and the sets that
are excluded or used
for sensitivity
analyses.}\tabularnewline
\toprule\noalign{}
\begin{minipage}[b]{\linewidth}\raggedright
Dataset
\end{minipage} &
\begin{minipage}[b]{\linewidth}\raggedright
Cases
\end{minipage} &
\begin{minipage}[b]{\linewidth}\raggedright
How it was built,
and what it is used
for
\end{minipage} \\
\midrule\noalign{}
\endfirsthead
\toprule\noalign{}
\begin{minipage}[b]{\linewidth}\raggedright
Dataset
\end{minipage} &
\begin{minipage}[b]{\linewidth}\raggedright
Cases
\end{minipage} &
\begin{minipage}[b]{\linewidth}\raggedright
How it was built,
and what it is used
for
\end{minipage} \\
\midrule\noalign{}
\endhead
\bottomrule\noalign{}
\endlastfoot
M-CSA source
\cite{ribeiro2018} &
434 curated entries
(release
2026-07-31), 2,563
catalytic residues &
the curated
mechanisms and
catalytic residues
of the Mechanism and
Catalytic Site
Atlas; structures
come from the
Protein Data Bank
\cite{burley2025,berman2000} \\
143-enzyme set
(primary) &
\textbf{143
enzymes}, 77 enzyme
sub-subclasses,
median 6 catalytic
residues, 2 with a
bound ligand & at
most 25 per
top-level EC class
(25, 25, 25, 25, 25
and 18 enzymes), at
least 2 resolved
catalytic residues,
ranked within a
class by quality,
catalytic-residue
count and
resolution;
packing-matched,
count-gated
controls \\
Pilot set & 53
enzymes (35
peptidases, 34 with
a ligand), 52 added
to the 143 & active
sites from UniProt
\cite{uniprot2023}
and M-CSA; one
enzyme shares PDB
entry 1LJL with a
main enzyme and is
counted once;
controls matched on
residue type and
burial only \\
\textbf{Structure-space
evaluation set} &
\textbf{195 enzymes}
(143 + 52) & the 29
structure-space
main-set metrics;
the main set alone
is kept beside the
pooled score (Table
S28) \\
\textbf{Prediction-based
evaluation set} &
\textbf{59 natural
enzymes, 286
catalytic/control
pairs} & the natural
arm of the
re-predicted ladder
(the 59 natural
enzymes of 2.3, 28
of which are in the
143); the two
reference rows pLDDT
and pTM are scored
on the same pairs \\
Excluded from
pooling & 30 de novo
designs; 15 designed
systems of the
ladder (74 systems,
312 pairs with the
59 natural ones) &
the control of a
design is
constructed, so
designs are not
pooled (Tables S29
and S30) \\
Shortened-motif set
& the 143 enzymes
with at most 3
catalytic residues
kept & a sensitivity
analysis of the
specificity ratio
(S31) \\
\end{longtable}

\textbf{Why this
dataset.} A lesion
at a catalytic
residue is
meaningful only for
enzymes whose
catalytic residues
are curated, and
M-CSA is the curated
source. The
143-enzyme set was
expanded from the
53-enzyme pilot for
EC balance and not
for N: the pilot was
66\% peptidase, and
with a flat result
more systems buy
precision on zero
while broader
coverage buys
generality. The
chemistry ladder
isolates chemical
identity at
near-constant
geometry, and its
isosteric step is
the hardest
synthetic negative;
the matched control,
which is mandatory
at every level,
tests whether a
metric knows that
the lesion is at the
catalytic site and
not merely that a
buried residue was
changed. The pilot
enzymes are pooled
because the
evaluation of a
ranking needs as
many enzymes as
possible in each
test; the effect of
pooling is reported
(adding the 52 pilot
enzymes changes the
mean R of a metric
by at most 0.06,
median 0.02, Table
S28). The
prediction-based arm
uses the natural
enzymes of the
substrate-swap set,
which have a cognate
ligand, because the
interface metrics
are undefined
without one.

\textbf{Metrics
scored.} The 29
structure-space
metrics (Table 1)
are computed on
every structure. The
four
prediction-based
main-set metrics
(Table 1, defined in
2.3) are computed on
each re-predicted
complex. The seed
noise of the ladder
is estimated from a
replicate with a
second seed on 30
natural pairs (Table
S32). The 35 ranked
items of Table 10
are the 29
structure-space and
4 prediction-based
main-set metrics and
the 2 reference rows
(pLDDT, pTM); the 8
supplementary
comparators and the
per-step flags are
in Table S28.

\textbf{Ranking
criteria.} The
\textbf{rank
statistic R} is the
probability form of
the comparison
between the two
arms. It is the
quantity that a
shortlist depends
on, because
pipelines rank
candidates and do
not use magnitudes;
Results 3.4 uses it
to rank the metrics
for each lesion. For
an item \(m\), a
step \(s\) and the
\(n_s\) enzymes that
carry both arms and
in which at least
one arm moves,

\begin{equation*}R_{m,s} \;=\; \frac{1}{n_s}\sum_{i=1}^{n_s}\Big[\mathbf{1}\big(\lvert\Delta^{\mathrm{cat}}_{i}\rvert > \lvert\Delta^{\mathrm{ctl}}_{i}\rvert\big) + \tfrac{1}{2}\,\mathbf{1}\big(\lvert\Delta^{\mathrm{cat}}_{i}\rvert = \lvert\Delta^{\mathrm{ctl}}_{i}\rvert\big)\Big] \tag{19}\end{equation*}

(an enzyme in which
both changes are
exactly zero carries
no ranking
information and is
dropped; a step
needs at least five
enzymes).
\(R = 0.5\) means no
specificity,
\(R > 0.5\) that the
item moves more at
the catalytic
lesion, \(R < 0.5\)
that it moves more
at the control;
unlike a
direction-free
AUROC, it has a
floor of 0 and not
of 0.5.

\emph{Which steps
count.} Where the
control arm does not
move beyond the
item's own noise,
\(R\) measures
whether the item
moves at all and not
whether it is
specific, and where
the catalytic
response is not
distinguishable from
zero there is
nothing to rank. A
structure-space step
therefore enters the
set \(S_m\) of
counted steps only
if the item responds
at that step (the
interval of its
median change
excludes zero) and
its control response
is above its noise,
as established on
the 143 main
enzymes. The
prediction-based
cells carry no such
flag, because the
control response of
the predictor is
noise-limited and is
read through the
distance-matched
subset instead; they
all count. The score
of an item is the
mean over the
counted steps,

\begin{equation*}\bar R_m \;=\; \frac{1}{\lvert S_m\rvert}\sum_{s \in S_m} R_{m,s} , \tag{20}\end{equation*}

and an item with no
counted step has no
mean. The number of
steps counted for
each item,
\(\lvert S_m\rvert\)
out of the steps
with a value, is in
Table S28.

\emph{Interval.} A
95\% percentile
bootstrap interval
of \(\bar R_m\) (Eq.
3 with \(\bar R_m\)
in place of
\(A^{*}\); 5,000
draws) that
resamples enzyme
sub-subclasses, the
first three EC
fields, a stated
proxy for
superfamily, and for
the predictor the
systems; the same
drawn clusters serve
every step of an
item.

\emph{Chance
reference.} As in
2.2, the top of a
ranking of \(M\)
items looks good by
chance. The labels
catalytic and
control are swapped
at random within
each enzyme
(\(s_i^{(r)} \sim \mathrm{Bernoulli}(1/2)\),
\(N_{\mathrm{perm}} = 5{,}000\)
replicates, the same
swaps for every
item),
\(\bar R^{(r)}_m\)
is recomputed for
every ranked item,
and
\(T^{(r)} = \max_m \bar R^{(r)}_m\)
(Eq. 4) has the
distribution of the
best score of \(M\)
unrelated items; its
median and 95th
percentile are the
\textbf{best-of-33
threshold}
(\(M = 33\) items
that have a counted
step; median 0.57,
95th percentile
0.62), a reference
level for reading
the ranking
\cite{westfall1993}.
A mean above it is
in bold in Table 10.

\emph{Onsets and
grouping.} The
metrics are also
differentiated by
the lesion they can
detect. The steps
are ordered
isosteric
\textless{}
non-isosteric
\textless{} Ala
\textless{} Gly
(15.3, 36.4, 53.2
and 80.5 ų). The
\textbf{onset of
detection} is the
smallest step at
which the item
responds (Eq. 18;
for the
structure-space
metrics the response
analysis of the 143
main enzymes, 2,000
draws, seed
20260807; for the
prediction-based
metrics the signed
changes of the
natural pairs,
bootstrap over
systems),

\begin{equation*}s_{\mathrm{det}}(m) \;=\; \min\{\,s : m \text{ responds at } s\,\}, \tag{21}\end{equation*}

with ``Never'' if
there is none. Rows
of Table 10 are
grouped by
\(s_{\mathrm{det}}\)
(the first column),
ranked by
\(\bar R_m\) within
each group, and a
metric with no
counted step closes
its group. The
\textbf{onset of
specificity} is the
smallest step at
which the item is
specific under the
specificity-ratio
rule
(structure-space
metrics only;
reported in Table
S28); the ratio and
the rule are

\begin{equation*}\mathrm{SR}_{m,s} \;=\; \frac{\operatorname{median}_i \lvert\Delta^{\mathrm{cat}}_{i}\rvert}{\operatorname{median}_i \lvert\Delta^{\mathrm{ctl}}_{i}\rvert}, \tag{22}\end{equation*}

and an item is
called
\textbf{specific}
only when all five
conditions hold: (1)
the catalytic
response is
significant; (2) the
interval of SR
excludes 1; (3) the
control arm
contributes at least
half as many
observations as the
catalytic arm; (4)
the control response
exceeds 1\% of the
item's own native
standard deviation,
because a ratio
whose denominator is
below the item's
noise is not
specificity; (5)
both arms changed
equal numbers of
residues (2,000
draws, seed
20260807; counts are
given as distinct
metrics, with
literal duplicates
collapsed). The
specificity ratio,
its interval and the
verdict (specific,
non-specific, blind
or invariant) of
every eligible
metric at every step
are in Table S33,
and those of the
whole-protein and
sequence-only
comparators in Table
S34; the 97 cells
that survive the
controls are listed
in Table S35; and
the per-step form of
R on the 143 main
enzymes alone, with
90\% intervals and
the comparators, is
in Table S36.

\emph{Distance-matched
R for the
predictor.}
Catalytic residues
touch the ligand by
definition, so an
interface item would
respond more to a
catalytic
substitution than to
a distant control
even if it knew
nothing about
catalysis. The
following plan was
fixed before any
distance was
computed. For every
pair the minimum
distance
\(d_{\mathrm{lig}}\)
from the substituted
side chain to the
ligand is measured,
and \(R\) of Eq.
(19) is also
computed on the
\textbf{caliper-matched
subset}: pairs with
the same residue
types, the control
not adjacent to a
catalytic residue,
and

\begin{equation*}\big\lvert d^{\mathrm{cat}}_{\mathrm{lig}} - d^{\mathrm{ctl}}_{\mathrm{lig}}\big\rvert \le 2\ \text{Å}, \qquad \big\lvert \mathrm{relSASA}^{\mathrm{cat}} - \mathrm{relSASA}^{\mathrm{ctl}}\big\rvert \le 0.10, \tag{23}\end{equation*}

extended by 54 pairs
given a new
in-caliper control
(53 new predictions)
wherever one
existed, pooled over
the four steps (68
to 69 pairs in 21 to
22 systems in Table
10). For the
original control the
median distance to
the ligand was 13.4
Å against 3.1 Å for
the catalytic
residue. The
distance of the
substituted and the
control residue to
the ligand for every
pair is in Table
S37, the candidate
controls for the
pairs that had none
in the caliper, and
why one was or was
not selected, in
Table S38, and the
ratio in thirds of
the distance to the
ligand in Table S39.

\emph{The
specificity-ratio
analysis of the
predictor (Table
S30).} The earlier
analysis reports the
ratio of the mean
change at the
catalytic residue to
the mean change at
the control (a
median-based form is
also given), with a
bootstrap over
systems (5,000
draws, seed
20260806), under
three successively
stricter controls:
(i) the original
control, on the
isosteric step only;
(ii) the
caliper-matched
subset of Eq. (23),
with all four steps
pooled (Table S40);
(iii) a regression
adjustment on all
pairs and steps
(Table S41),

\begin{equation*}\mathrm{SR} = \frac{\overline{\lvert\Delta\rvert}_{\mathrm{cat}}}{\overline{\lvert\Delta\rvert}_{\mathrm{ctl}}}, \qquad \log\big(\lvert\Delta\rvert + \varepsilon\big) = \beta_0 + \beta_{\mathrm{arm}}\,\mathbf{1}(\text{catalytic}) + \beta_d\, d_{\mathrm{lig}} + \beta_s\, \mathrm{relSASA} + \text{error}, \tag{24}\end{equation*}

where
\(\varepsilon\) is a
small
metric-specific
constant that keeps
zero changes finite
and the adjusted
ratio is
\(\exp(\beta_{\mathrm{arm}})\),
with a bootstrap
over systems. The
\textbf{decision
rule for the
predictor} was fixed
before the analysis:
an interface item
counts as specific
only if the lower
90\% bound of its
mean-based SR in the
matched subset
exceeds 1.2, with at
least 30 pairs in at
least 15 systems;
the active-site
accuracy is reported
as knock-on
displacement with no
specificity verdict,
because the
substituted residue
is excluded from its
own comparison. The
whole
structure-space
panel is summarised
in the Supplement by
the geometric-mean
SR over the metrics
with a defined
ratio, with a nested
bootstrap (3,000
draws) that
recomputes every
metric inside each
draw, so that
correlation between
metrics is carried
through; equivalence
to 1 is judged
against a ±25\% band
(S42). The
re-predicted ladder
at each step, for
all metrics, natural
and de novo, is in
Table S43.

\subsubsection{2.5.2
Detecting second
shell lesion}

\textbf{Dataset.}
The 143 main enzymes
of 2.5.1 pooled with
the 52 pilot enzymes
that are not in the
main set, with a
control arm for up
to 193 of them;
structure-space
metrics are scored
on this lesion. The
number of enzymes in
which an item moves
at all, and so
carries information,
ranges from 0 to 182
per item and dose.

\textbf{Why this
dataset.} The
second-shell lesion
asks whether the
metrics see past the
first shell:
second-shell
residues around the
catalytic residues
are removed in
graded doses of 1, 2
and 4 residues, and
the control removes
the same number of
residues around a
matched
non-catalytic site.
The same enzymes as
the catalytic lesion
are used, so that
the two lesions can
be read on one panel
and one population,
and the same pooling
applies.

\textbf{Metrics
scored.} The 29
structure-space
main-set metrics,
scored over the
catalytic residues
in both arms (so
that a site-scoped
item never knows
which residues are
catalytic). The 8
supplementary
comparators and the
per-dose flags are
in Table S44.

\textbf{Ranking
criteria.} The doses
are
\(d \in \{1, 2, 4\}\)
residues removed in
place of the steps
of Eq. (19), and the
rank statistic is

\begin{equation*}R_{m,d} \;=\; \frac{1}{n_d}\sum_{i=1}^{n_d}\Big[\mathbf{1}\big(\lvert\Delta^{\mathrm{cat}}_{i}\rvert > \lvert\Delta^{\mathrm{ctl}}_{i}\rvert\big) + \tfrac{1}{2}\,\mathbf{1}\big(\lvert\Delta^{\mathrm{cat}}_{i}\rvert = \lvert\Delta^{\mathrm{ctl}}_{i}\rvert\big)\Big], \qquad \bar R_m = \frac{1}{3}\sum_{d} R_{m,d} . \tag{25}\end{equation*}

with the same
bootstrap over
enzyme
sub-subclasses, the
same onset of
detection (Eq. 21,
over doses) and the
same grouping as for
the catalytic
lesion. The doses
with a response and
with a control
response above the
item's noise are
counted for each
item, and a dose
enters a counted set
only where both
hold, as in 2.5.1;
here no cell
qualifies for any
item, so the ranking
uses the raw
\(R_{m,d}\) of every
dose with at least
five informative
enzymes, raw values
are reported, and no
best-of-\(M\)
threshold is used.
Table 11 ranks the
items within each
detection group by
raw \(\bar R_m\) and
gives, for each
item, the number of
doses (of 3) at
which it responds
and at which the
control response is
above its noise.

\begin{center}\rule{0.5\linewidth}{0.5pt}\end{center}

\textbf{Provenance
and data
availability.} Every
number in this paper
is read from a table
generated by script
from committed
result files, with
seeds and input
hashes recorded.
Each experiment's
analysis plan was
committed and dated
before the analysis
it governs, and
amendments are
listed with the date
they were made. The
Supplement gives,
for every metric and
every experiment,
whether the metric
was applicable and
why; the full
per-metric tables
behind Tables 5--11;
and the experiments
that could not be
made. All tables are
released as CSV
files (main-text
tables in
\texttt{main/},
supplementary tables
in
\texttt{supplementary/});
the Supplementary
PDF describes every
table and gives the
path of its file.
Structures and
predictor weights
are third-party and
are not
redistributed; their
hashes are. BglB
variant-level data
are not
redistributed
because their
licence is unstated;
only aggregates are
reported.

\section{3 Results}

Each subsection
follows one order:
the experiment, the
dataset, the
criteria by which
the metrics are
ranked, and the
result table. Counts
of metrics are of
distinct metrics in
the main set of 33
(29 structure-space,
4 prediction-based;
Table 1 of the
Methods), and
intervals are 95\%
bootstrap intervals
that resample the
unit named in each
subsection, unless
stated otherwise.
Numbers in square
brackets are
references, listed
in the References.
Complete per-metric
tables and further
analyses are in the
supplementary tables
(S numbers), which
are described in the
Supplementary PDF
and released as CSV
files
(\texttt{main/} and
\texttt{supplementary/}).

\begin{center}\rule{0.5\linewidth}{0.5pt}\end{center}

\subsection{3.1
Activity
discrimination test}

\textbf{Experiment.}
Metrics are scored
on real proteins in
which one state is
catalytically dead
and the other
active, and each
metric is ranked by
how well its value
separates the two
states (Methods
2.2). Two datasets
are used, each with
its own table.

\subsubsection{3.1.1
Discriminate natural
active and dead
enzymes}

\textbf{Dataset.} 21
zymogen--mature
pairs from 21
peptidases (9
serine, 8 cysteine,
4 aspartic),
downloaded from the
RCSB Protein Data
Bank
\cite{burley2025,berman2000}.
The zymogen is the
dead form and the
mature enzyme the
active form, and the
two forms share a
ligand. Catalytic
residues come from
M-CSA
\cite{ribeiro2018}
and UniProt
\cite{uniprot2023}
annotation;
prediction-based
quantities are
Chai-1
\cite{chai2024}
predictions of both
forms with the
shared ligand. Every
item is scored on
all 21 pairs
(Methods 2.2.1).

\textbf{Ranking
criteria.} The 33
main-set metrics and
6 reference rows are
ranked by the
direction-free AUROC
\cite{hanley1982,mann1947}
of dead against
active (0.5 = no
separation), with a
95\% bootstrap
interval that
resamples pairs.
``Higher in'' is the
form in which the
metric takes larger
values. An AUROC in
bold exceeds the
best-of-39 threshold
\cite{westfall1993},
the 95th percentile
of the best AUROC
among 39 unrelated
items under random
label swaps within
pairs (median 0.63,
95th percentile
0.68).

\begin{longtable}[]{@{}
  >{\raggedright\arraybackslash}p{(\columnwidth - 10\tabcolsep) * \real{0.0709}}
  >{\raggedright\arraybackslash}p{(\columnwidth - 10\tabcolsep) * \real{0.2992}}
  >{\raggedright\arraybackslash}p{(\columnwidth - 10\tabcolsep) * \real{0.1890}}
  >{\raggedright\arraybackslash}p{(\columnwidth - 10\tabcolsep) * \real{0.1260}}
  >{\raggedright\arraybackslash}p{(\columnwidth - 10\tabcolsep) * \real{0.2047}}
  >{\raggedright\arraybackslash}p{(\columnwidth - 10\tabcolsep) * \real{0.1102}}@{}}
\caption*{\textbf{Table
5. Separation of
real dead from
active enzymes on
the 21-pair
evaluation set (top
30 of 39 ranked
items).} The 39
items are the 33
main-set metrics and
the 6 reference rows
(italic), ranked by
the direction-free
AUROC of dead
against active
(floor 0.5) over 21
pairs from 21
proteins; every row
has n = 21 pairs.
Intervals are 95\%
bootstrap intervals
that resample pairs.
``Higher in'' is the
form in which the
metric takes larger
values. An AUROC in
bold exceeds the
best-of-39
threshold, the 95th
percentile of the
best AUROC among 39
unrelated metrics in
label-swap
simulations (median
0.63, 95th
percentile 0.68;
defined in the
Methods; a reference
level for reading
the ranking). The
nine lowest-ranked
items (AUROC 0.50 to
0.52) and all 158
metrics scored on
these pairs are in
Table S11. Source:
main/T5\_dead\_vs\_active.csv.}\tabularnewline
\toprule\noalign{}
\begin{minipage}[b]{\linewidth}\raggedright
Rank
\end{minipage} &
\begin{minipage}[b]{\linewidth}\raggedright
Item
\end{minipage} &
\begin{minipage}[b]{\linewidth}\raggedright
Group
\end{minipage} &
\begin{minipage}[b]{\linewidth}\raggedright
Mode
\end{minipage} &
\begin{minipage}[b]{\linewidth}\raggedright
AUROC {[}95\%
interval{]}
\end{minipage} &
\begin{minipage}[b]{\linewidth}\raggedright
Higher in
\end{minipage} \\
\midrule\noalign{}
\endfirsthead
\toprule\noalign{}
\begin{minipage}[b]{\linewidth}\raggedright
Rank
\end{minipage} &
\begin{minipage}[b]{\linewidth}\raggedright
Item
\end{minipage} &
\begin{minipage}[b]{\linewidth}\raggedright
Group
\end{minipage} &
\begin{minipage}[b]{\linewidth}\raggedright
Mode
\end{minipage} &
\begin{minipage}[b]{\linewidth}\raggedright
AUROC {[}95\%
interval{]}
\end{minipage} &
\begin{minipage}[b]{\linewidth}\raggedright
Higher in
\end{minipage} \\
\midrule\noalign{}
\endhead
\bottomrule\noalign{}
\endlastfoot
1 & \emph{pLDDT}
\cite{jumper2021} &
Reference row &
prediction &
\textbf{0.79
{[}0.68, 0.90{]}} &
active \\
2 & \emph{pTM}
\cite{jumper2021} &
Reference row &
prediction &
\textbf{0.75
{[}0.62, 0.87{]}} &
active \\
3 & Intra-residue
repulsion (Rosetta,
site)
\cite{alford2017} &
Rosetta energy &
structure & 0.61
{[}0.52, 0.72{]} &
dead \\
4 & Repulsion
(Rosetta, site)
\cite{alford2017} &
Rosetta energy &
structure & 0.60
{[}0.51, 0.72{]} &
dead \\
5 &
\emph{Crystallographic
resolution} &
Reference row &
structure & 0.58
{[}0.50, 0.72{]} &
dead \\
6 & Fraction buried
& Geometry &
structure & 0.57
{[}0.50, 0.69{]} &
dead \\
7 & Attraction
(Rosetta, site)
\cite{alford2017} &
Rosetta energy &
structure & 0.56
{[}0.51, 0.64{]} &
active \\
8 & lk-ball
solvation (Rosetta,
site)
\cite{alford2017} &
Rosetta energy &
structure & 0.56
{[}0.50, 0.66{]} &
dead \\
9 & Ramachandran
(Rosetta, site)
\cite{alford2017} &
Rosetta energy &
structure & 0.56
{[}0.51, 0.63{]} &
dead \\
10 & \emph{Sequence
length} & Reference
row & structure &
0.56 {[}0.50,
0.65{]} & dead \\
11 & Relative SASA,
mean & Geometry &
structure & 0.55
{[}0.50, 0.66{]} &
active \\
12 & Total SASA &
Geometry & structure
& 0.55 {[}0.50,
0.67{]} & active \\
13 & Rotamer,
Dunbrack (Rosetta,
site)
\cite{alford2017} &
Rosetta energy &
structure & 0.55
{[}0.50, 0.66{]} &
dead \\
14 & Pairwise
distance, max &
Geometry & structure
& 0.54 {[}0.50,
0.64{]} & dead \\
15 & ipTM (Chai-1)
\cite{chai2024,evans2021}
& Interface
confidence &
prediction & 0.54
{[}0.50, 0.61{]} &
active \\
16 & pKa, range
(PROPKA)
\cite{olsson2011,sondergaard2011}
& pKa & structure &
0.54 {[}0.50,
0.69{]} & dead \\
17 & Polar contacts
& Geometry &
structure & 0.54
{[}0.50, 0.59{]} &
dead \\
18 & Ligand
clearance
\cite{ahern2026} &
Active-site accuracy
& prediction & 0.54
{[}0.50, 0.65{]} &
dead \\
19 & Net van der
Waals (Rosetta,
site)
\cite{alford2017} &
Rosetta energy &
structure & 0.54
{[}0.50, 0.61{]} &
active \\
20 & Relative SASA,
max & Geometry &
structure & 0.53
{[}0.50, 0.65{]} &
active \\
21 & Radius of
gyration & Geometry
& structure & 0.53
{[}0.50, 0.64{]} &
dead \\
22 & AME RMSD
\cite{ahern2026} &
Active-site accuracy
& prediction & 0.53
{[}0.50, 0.66{]} &
active \\
23 & Pairwise
distance, mean &
Geometry & structure
& 0.53 {[}0.50,
0.64{]} & dead \\
24 & Solvation
(Rosetta, site)
\cite{alford2017} &
Rosetta energy &
structure & 0.53
{[}0.50, 0.60{]} &
dead \\
25 & Omega torsion
(Rosetta, site)
\cite{alford2017} &
Rosetta energy &
structure & 0.53
{[}0.50, 0.63{]} &
dead \\
26 & pKa, min
(PROPKA)
\cite{olsson2011,sondergaard2011}
& pKa & structure &
0.52 {[}0.50,
0.65{]} & active \\
27 & Amino-acid
propensity (Rosetta,
site)
\cite{alford2017} &
Rosetta energy &
structure & 0.52
{[}0.50, 0.59{]} &
active \\
28 & Electrostatics
(Rosetta, site)
\cite{alford2017} &
Rosetta energy &
structure & 0.52
{[}0.50, 0.58{]} &
active \\
29 & pKa, max
(PROPKA)
\cite{olsson2011,sondergaard2011}
& pKa & structure &
0.52 {[}0.50,
0.67{]} & dead \\
30 & H-bond,
backbone to side
chain (Rosetta,
site)
\cite{alford2017} &
Rosetta energy &
structure & 0.52
{[}0.50, 0.58{]} &
dead \\
\end{longtable}

\textbf{More
detail.} Every
metric scored on the
21 pairs, 158 in
all, with the
best-of-158
threshold: Table
S11. The wider set
of 49 pairs (136
metrics) with the
AUROC against the
same-state null:
Table S7. The
definitions of the
six pair sets and
where each is used:
Table S9. The 28
pairs re-predicted
with Chai-1,
restated with the
definedness gate:
Table S8. The check
of the
substrate-trapping
mutants against the
literature (26 of 30
not confirmed dead),
a tier that was
retired: Table S10.
The reference rows
in full, with
intervals: Table
S12.

\subsubsection{3.1.2
Discriminate de novo
designed active and
no active enzymes}

\textbf{Dataset.}
192 designs from a
deposited
metallohydrolase
design campaign
\cite{kim2026}: 16
have a reported
kcat/KM
(\emph{active}) and
176 are \emph{no
active}. The
comparison is active
against no active
designs. The designs
lie in 136 backbone
clusters, the 16
active designs in 8
of them (10 of the
16 in two clusters).
Structure-space
metrics are computed
on the designs as
deposited and
prediction-based
quantities on Chai-1
\cite{chai2024}
predictions with the
design's
transition-state-analogue
ligand and zinc.
Every item is scored
on all 192 designs
(Methods 2.2.2).

\textbf{Ranking
criteria.} The 36
items are 29
structure-space and
2 prediction-based
main-set metrics and
5 reference rows
(tyrosine fraction,
net charge, sequence
length, pLDDT, pTM).
They are ranked by
the direction-free
AUROC
\cite{hanley1982,mann1947}
of active against no
active designs (0.5
= no separation),
with a 95\%
bootstrap interval
that resamples the
136 backbone
clusters
\cite{field2007}.
``Higher in'' is the
group in which the
metric takes larger
values. An AUROC in
bold exceeds the
best-of-36 threshold
\cite{westfall1993},
the 95th percentile
of the best AUROC
among 36 items when
the activity labels
are permuted among
designs (median
0.66, 95th
percentile 0.73);
the permutation
treats designs as
exchangeable, which
understates the
chance range when
the actives are
clumped in few
clusters. AME is
measured on the
plates against each
design's own model
(a different metric
from AME against a
deposited structure)
and is given in
Table S13.

\begin{longtable}[]{@{}
  >{\raggedright\arraybackslash}p{(\columnwidth - 10\tabcolsep) * \real{0.0709}}
  >{\raggedright\arraybackslash}p{(\columnwidth - 10\tabcolsep) * \real{0.2598}}
  >{\raggedright\arraybackslash}p{(\columnwidth - 10\tabcolsep) * \real{0.1732}}
  >{\raggedright\arraybackslash}p{(\columnwidth - 10\tabcolsep) * \real{0.1181}}
  >{\raggedright\arraybackslash}p{(\columnwidth - 10\tabcolsep) * \real{0.1969}}
  >{\raggedright\arraybackslash}p{(\columnwidth - 10\tabcolsep) * \real{0.1811}}@{}}
\caption*{\textbf{Table
6. Separation of
active from no
active designs on
the 192-design
evaluation set (top
30 of 36 ranked
items).} The 36
items are 29
structure-space and
2 prediction-based
main-set metrics and
5 reference rows
(italic), ranked by
the direction-free
AUROC (floor 0.5) of
the 16 active
designs against the
176 no active
designs; every row
has n = 192 designs
(16 active).
Intervals are 95\%
bootstrap intervals
that resample the
136 backbone
clusters. ``Higher
in'' is the group in
which the metric
takes larger values:
``active'' (the 16
designs with a
reported kcat/KM) or
``no active'' (the
other 176 designs).
An AUROC in bold
exceeds the
best-of-36
threshold, the 95th
percentile of the
best AUROC among 36
unrelated items when
the activity labels
are permuted among
designs (median
0.66, 95th
percentile 0.73; a
reference level for
reading the
ranking). Items 31
to 36 and all 128
metrics scored on
the designs are in
Table S13. Source:
main/T6\_plated\_designs\_auroc.csv.}\tabularnewline
\toprule\noalign{}
\begin{minipage}[b]{\linewidth}\raggedright
Rank
\end{minipage} &
\begin{minipage}[b]{\linewidth}\raggedright
Item
\end{minipage} &
\begin{minipage}[b]{\linewidth}\raggedright
Group
\end{minipage} &
\begin{minipage}[b]{\linewidth}\raggedright
Mode
\end{minipage} &
\begin{minipage}[b]{\linewidth}\raggedright
AUROC {[}95\%
interval{]}
\end{minipage} &
\begin{minipage}[b]{\linewidth}\raggedright
Higher in
\end{minipage} \\
\midrule\noalign{}
\endfirsthead
\toprule\noalign{}
\begin{minipage}[b]{\linewidth}\raggedright
Rank
\end{minipage} &
\begin{minipage}[b]{\linewidth}\raggedright
Item
\end{minipage} &
\begin{minipage}[b]{\linewidth}\raggedright
Group
\end{minipage} &
\begin{minipage}[b]{\linewidth}\raggedright
Mode
\end{minipage} &
\begin{minipage}[b]{\linewidth}\raggedright
AUROC {[}95\%
interval{]}
\end{minipage} &
\begin{minipage}[b]{\linewidth}\raggedright
Higher in
\end{minipage} \\
\midrule\noalign{}
\endhead
\bottomrule\noalign{}
\endlastfoot
1 & Intra-residue
repulsion (Rosetta,
site)
\cite{alford2017} &
Rosetta energy &
structure &
\textbf{0.80
{[}0.64, 0.89{]}} &
no active \\
2 & Repulsion
(Rosetta, site)
\cite{alford2017} &
Rosetta energy &
structure &
\textbf{0.78
{[}0.57, 0.92{]}} &
no active \\
3 & \emph{Tyrosine
fraction} &
Reference row &
structure &
\textbf{0.75
{[}0.55, 0.86{]}} &
active \\
4 & \emph{pLDDT}
\cite{jumper2021} &
Reference row &
prediction & 0.73
{[}0.53, 0.86{]} &
active \\
5 & \emph{pTM}
\cite{jumper2021} &
Reference row &
prediction & 0.72
{[}0.54, 0.85{]} &
active \\
6 & Site total
(Rosetta)
\cite{alford2017} &
Rosetta energy &
structure & 0.71
{[}0.52, 0.82{]} &
no active \\
7 & Rotamer,
Dunbrack (Rosetta,
site)
\cite{alford2017} &
Rosetta energy &
structure & 0.71
{[}0.58, 0.88{]} &
no active \\
8 & Per-chain pTM
minimum (Chai-1)
\cite{chai2024,jumper2021}
& Interface
confidence &
prediction & 0.70
{[}0.54, 0.83{]} &
active \\
9 & ipTM (Chai-1)
\cite{chai2024,evans2021}
& Interface
confidence &
prediction & 0.70
{[}0.52, 0.84{]} &
active \\
10 & \emph{Sequence
length} & Reference
row & structure &
0.68 {[}0.51,
0.86{]} & no
active \\
11 & Worst residue
(Rosetta, site)
\cite{alford2017} &
Rosetta energy &
structure & 0.68
{[}0.52, 0.78{]} &
no active \\
12 & Attraction
(Rosetta, site)
\cite{alford2017} &
Rosetta energy &
structure & 0.67
{[}0.57, 0.82{]} &
active \\
13 & lk-ball
solvation (Rosetta,
site)
\cite{alford2017} &
Rosetta energy &
structure & 0.66
{[}0.54, 0.82{]} &
active \\
14 & Total SASA &
Geometry & structure
& 0.65 {[}0.54,
0.89{]} & active \\
15 & Relative SASA,
mean & Geometry &
structure & 0.65
{[}0.53, 0.89{]} &
active \\
16 & Electrostatics
(Rosetta, site)
\cite{alford2017} &
Rosetta energy &
structure & 0.64
{[}0.51, 0.79{]} &
no active \\
17 & Pairwise
distance, max &
Geometry & structure
& 0.64 {[}0.51,
0.77{]} & active \\
18 & Relative SASA,
max & Geometry &
structure & 0.62
{[}0.52, 0.84{]} &
active \\
19 & Pairwise
distance, mean &
Geometry & structure
& 0.61 {[}0.51,
0.77{]} & active \\
20 & Radius of
gyration & Geometry
& structure & 0.61
{[}0.51, 0.77{]} &
active \\
21 & H-bond, side
chain (Rosetta,
site)
\cite{alford2017} &
Rosetta energy &
structure & 0.61
{[}0.51, 0.73{]} &
no active \\
22 & Solvation
(Rosetta, site)
\cite{alford2017} &
Rosetta energy &
structure & 0.60
{[}0.51, 0.71{]} &
no active \\
23 & Amino-acid
propensity (Rosetta,
site)
\cite{alford2017} &
Rosetta energy &
structure & 0.60
{[}0.51, 0.73{]} &
active \\
24 & Polar contacts
& Geometry &
structure & 0.59
{[}0.50, 0.74{]} &
active \\
25 & Omega torsion
(Rosetta, site)
\cite{alford2017} &
Rosetta energy &
structure & 0.58
{[}0.50, 0.81{]} &
active \\
26 & \emph{Net
charge} & Reference
row & structure &
0.56 {[}0.51,
0.83{]} & active \\
27 & pKa, max
(PROPKA)
\cite{olsson2011,sondergaard2011}
& pKa & structure &
0.56 {[}0.50,
0.78{]} & no
active \\
28 & pKa, range
(PROPKA)
\cite{olsson2011,sondergaard2011}
& pKa & structure &
0.55 {[}0.50,
0.77{]} & no
active \\
29 & Net van der
Waals (Rosetta,
site)
\cite{alford2017} &
Rosetta energy &
structure & 0.54
{[}0.50, 0.80{]} &
active \\
30 & Fraction buried
& Geometry &
structure & 0.53
{[}0.50, 0.70{]} &
no active \\
\end{longtable}

\textbf{More
detail.} Items 31 to
36 and every other
metric scored on the
192 designs (128 in
all), with the
best-of-128
threshold: Table
S13. The reference
rows in full: Table
S12. The descriptive
ranking of the
metrics among the 16
active designs is in
3.3.2.

\begin{center}\rule{0.5\linewidth}{0.5pt}\end{center}

\subsection{3.2
Substrate
discrimination test}

\textbf{Experiment.}
Each enzyme is
predicted with
Chai-1
\cite{chai2024} in
four conditions that
differ only in the
ligand supplied: its
cognate ligand, a
wrong ligand of the
same enzyme class, a
wrong ligand of a
different class, and
a decoy ligand taken
from an unrelated
structure. Each
prediction-based
quantity is ranked
by how well it
separates the
cognate prediction
from a wrong-ligand
prediction of the
same enzyme (Methods
2.3).

\textbf{Dataset.} 55
natural enzymes of
M-CSA
\cite{ribeiro2018},
of the 59 in the
experiment: on
these, every scored
quantity has a
defined value in all
four conditions and
every supplied
ligand was present
in the prediction
(four enzymes are
excluded because a
heme-type ligand
could not be built
into a chemical
structure and the
predictor dropped
it; the analysis on
all 59 is in Table
S14). One n, 55
enzymes, applies to
every row (Methods
2.3).

\textbf{Ranking
criteria.} Chai-1
returns five models
per prediction, and
each score is
summarised over them
as a mean, a maximum
or a spread
(standard deviation
across models), and
for the AME RMSD
also as a minimum,
so one metric
appears in several
rows. The 30 scored
quantities are the
four
prediction-based
main-set metrics
(ipTM, per-chain pTM
minimum, ligand
clearance, AME
RMSD),
whole-prediction
confidence scores
(including the
reference rows pLDDT
and pTM), the AME
pass bit and four
bookkeeping
counters. Each is
ranked by the mean
of three
direction-free
AUROCs
\cite{hanley1982,mann1947},
one per kind of
wrong ligand (same
class, different
class, decoy), each
with a 95\% interval
that resamples
enzymes. ``Higher
in'' is the
condition in which
the quantity takes
larger values. A
mean in bold exceeds
the best-of-30
threshold
\cite{westfall1993},
the 95th percentile
of the best AUROC
among 30 quantities
when the four ligand
conditions are
permuted within each
enzyme (median 0.56,
95th percentile
0.61).

\begin{longtable}[]{@{}
  >{\raggedright\arraybackslash}p{(\columnwidth - 16\tabcolsep) * \real{0.0573}}
  >{\raggedright\arraybackslash}p{(\columnwidth - 16\tabcolsep) * \real{0.1656}}
  >{\raggedright\arraybackslash}p{(\columnwidth - 16\tabcolsep) * \real{0.1019}}
  >{\raggedright\arraybackslash}p{(\columnwidth - 16\tabcolsep) * \real{0.1529}}
  >{\raggedright\arraybackslash}p{(\columnwidth - 16\tabcolsep) * \real{0.1274}}
  >{\raggedright\arraybackslash}p{(\columnwidth - 16\tabcolsep) * \real{0.1274}}
  >{\raggedright\arraybackslash}p{(\columnwidth - 16\tabcolsep) * \real{0.1274}}
  >{\raggedright\arraybackslash}p{(\columnwidth - 16\tabcolsep) * \real{0.0573}}
  >{\raggedright\arraybackslash}p{(\columnwidth - 16\tabcolsep) * \real{0.0828}}@{}}
\caption*{\textbf{Table
7. The substrate
swap on the
55-enzyme evaluation
set, 30 scored
prediction-based
quantities ranked by
AUROC.}
Direction-free AUROC
(floor 0.5) of the
cognate prediction
against a
wrong-ligand
prediction of the
same enzyme, in
three columns by
kind of wrong
ligand, with 95\%
intervals that
resample enzymes;
the value of an
enzyme in a
condition is the
mean over seeds 101,
202 and 303. Rows
are ranked by the
mean of the three
AUROCs; a mean in
bold exceeds the
best-of-30 threshold
(median 0.56, 95th
percentile 0.61;
permutation of the
four ligand
conditions within
each enzyme, defined
in the Methods; a
reference level for
reading the
ranking). Chai-1
predicts five models
for every
prediction, and each
score is summarised
over them (mean,
maximum, spread as
the standard
deviation across
models, and for the
AME RMSD also the
minimum), so one
metric appears in
several rows, one
per summary.
Reference rows are
in italic. Source:
main/T7\_substrate\_swap.csv.}\tabularnewline
\toprule\noalign{}
\begin{minipage}[b]{\linewidth}\raggedright
Rank
\end{minipage} &
\begin{minipage}[b]{\linewidth}\raggedright
Quantity
\end{minipage} &
\begin{minipage}[b]{\linewidth}\raggedright
Summary over the
five models
\end{minipage} &
\begin{minipage}[b]{\linewidth}\raggedright
Group
\end{minipage} &
\begin{minipage}[b]{\linewidth}\raggedright
Same class
\end{minipage} &
\begin{minipage}[b]{\linewidth}\raggedright
Different class
\end{minipage} &
\begin{minipage}[b]{\linewidth}\raggedright
Decoy
\end{minipage} &
\begin{minipage}[b]{\linewidth}\raggedright
Mean
\end{minipage} &
\begin{minipage}[b]{\linewidth}\raggedright
Higher in
\end{minipage} \\
\midrule\noalign{}
\endfirsthead
\toprule\noalign{}
\begin{minipage}[b]{\linewidth}\raggedright
Rank
\end{minipage} &
\begin{minipage}[b]{\linewidth}\raggedright
Quantity
\end{minipage} &
\begin{minipage}[b]{\linewidth}\raggedright
Summary over the
five models
\end{minipage} &
\begin{minipage}[b]{\linewidth}\raggedright
Group
\end{minipage} &
\begin{minipage}[b]{\linewidth}\raggedright
Same class
\end{minipage} &
\begin{minipage}[b]{\linewidth}\raggedright
Different class
\end{minipage} &
\begin{minipage}[b]{\linewidth}\raggedright
Decoy
\end{minipage} &
\begin{minipage}[b]{\linewidth}\raggedright
Mean
\end{minipage} &
\begin{minipage}[b]{\linewidth}\raggedright
Higher in
\end{minipage} \\
\midrule\noalign{}
\endhead
\bottomrule\noalign{}
\endlastfoot
1 & ipTM
\cite{chai2024,evans2021}
& mean & Interface
confidence (main
set) & 0.69 {[}0.60,
0.78{]} & 0.75
{[}0.66, 0.84{]} &
0.65 {[}0.55,
0.76{]} &
\textbf{0.70} &
cognate \\
2 & ipTM
\cite{chai2024,evans2021}
& maximum &
Interface confidence
(main set) & 0.68
{[}0.59, 0.77{]} &
0.74 {[}0.65,
0.83{]} & 0.65
{[}0.54, 0.74{]} &
\textbf{0.69} &
cognate \\
3 & ipTM
\cite{chai2024,evans2021}
& spread (SD) &
Interface confidence
(main set) & 0.70
{[}0.60, 0.79{]} &
0.69 {[}0.60,
0.77{]} & 0.67
{[}0.59, 0.75{]} &
\textbf{0.69} &
wrong ligand \\
4 & pTM
\cite{jumper2021} &
spread (SD) &
Whole-prediction
confidence & 0.62
{[}0.55, 0.70{]} &
0.64 {[}0.58,
0.71{]} & 0.64
{[}0.57, 0.71{]} &
\textbf{0.64} &
wrong ligand \\
5 & Ligand clearance
\cite{ahern2026} &
best model &
Active-site accuracy
(main set) & 0.61
{[}0.52, 0.69{]} &
0.65 {[}0.57,
0.74{]} & 0.62
{[}0.53, 0.71{]} &
\textbf{0.63} &
wrong ligand \\
6 & \emph{pTM}
\cite{jumper2021} &
mean & Reference row
& 0.60 {[}0.56,
0.66{]} & 0.63
{[}0.58, 0.69{]} &
0.64 {[}0.58,
0.70{]} &
\textbf{0.62} &
cognate \\
7 & pTM
\cite{jumper2021} &
maximum &
Whole-prediction
confidence & 0.60
{[}0.56, 0.66{]} &
0.63 {[}0.58,
0.69{]} & 0.63
{[}0.58, 0.70{]} &
\textbf{0.62} &
cognate \\
8 & Per-chain pTM
minimum
\cite{chai2024,jumper2021}
& spread (SD) &
Interface confidence
(main set) & 0.58
{[}0.51, 0.68{]} &
0.58 {[}0.51,
0.68{]} & 0.69
{[}0.59, 0.79{]} &
\textbf{0.62} &
wrong ligand \\
9 & AME RMSD
\cite{ahern2026} &
spread (SD) &
Active-site accuracy
(main set) & 0.55
{[}0.51, 0.60{]} &
0.56 {[}0.51,
0.62{]} & 0.57
{[}0.52, 0.63{]} &
0.56 & wrong
ligand \\
10 & Per-chain pTM
minimum
\cite{chai2024,jumper2021}
& mean & Interface
confidence (main
set) & 0.58 {[}0.51,
0.67{]} & 0.60
{[}0.51, 0.71{]} &
0.50 {[}0.50,
0.62{]} & 0.56 &
cognate \\
11 & Per-chain pTM
minimum
\cite{chai2024,jumper2021}
& maximum &
Interface confidence
(main set) & 0.58
{[}0.51, 0.67{]} &
0.59 {[}0.51,
0.71{]} & 0.50
{[}0.50, 0.62{]} &
0.56 & cognate \\
12 & AME RMSD
\cite{ahern2026} &
maximum &
Active-site accuracy
(main set) & 0.53
{[}0.50, 0.57{]} &
0.54 {[}0.51,
0.58{]} & 0.54
{[}0.51, 0.59{]} &
0.54 & wrong
ligand \\
13 & AME RMSD
\cite{ahern2026} &
minimum &
Active-site accuracy
(main set) & 0.53
{[}0.50, 0.56{]} &
0.54 {[}0.51,
0.57{]} & 0.52
{[}0.50, 0.57{]} &
0.53 & wrong
ligand \\
14 & AME RMSD, no
symmetry resolution
\cite{ahern2026} &
minimum &
Active-site accuracy
(main set) & 0.52
{[}0.50, 0.56{]} &
0.54 {[}0.51,
0.57{]} & 0.53
{[}0.50, 0.57{]} &
0.53 & wrong
ligand \\
15 & AME RMSD
\cite{ahern2026} &
mean & Active-site
accuracy (main set)
& 0.53 {[}0.50,
0.56{]} & 0.53
{[}0.51, 0.57{]} &
0.52 {[}0.50,
0.57{]} & 0.53 &
wrong ligand \\
16 & pLDDT
\cite{jumper2021} &
spread (SD) &
Whole-prediction
confidence & 0.51
{[}0.50, 0.55{]} &
0.52 {[}0.50,
0.57{]} & 0.54
{[}0.50, 0.60{]} &
0.52 & mixed \\
17 & Per-chain pTM
of the protein
\cite{chai2024,jumper2021}
& mean &
Whole-prediction
confidence & 0.52
{[}0.50, 0.54{]} &
0.52 {[}0.50,
0.54{]} & 0.53
{[}0.51, 0.56{]} &
0.52 & cognate \\
18 & Per-chain pTM
of the protein
\cite{chai2024,jumper2021}
& maximum &
Whole-prediction
confidence & 0.52
{[}0.50, 0.54{]} &
0.52 {[}0.50,
0.55{]} & 0.53
{[}0.51, 0.55{]} &
0.52 & cognate \\
19 & pLDDT
\cite{jumper2021} &
best model &
Whole-prediction
confidence & 0.52
{[}0.50, 0.54{]} &
0.53 {[}0.51,
0.55{]} & 0.51
{[}0.50, 0.54{]} &
0.52 & cognate \\
20 & \emph{pLDDT}
\cite{jumper2021} &
mean & Reference row
& 0.52 {[}0.50,
0.54{]} & 0.53
{[}0.51, 0.55{]} &
0.51 {[}0.50,
0.54{]} & 0.52 &
cognate \\
21 & Per-chain pTM
of the protein
\cite{chai2024,jumper2021}
& spread (SD) &
Whole-prediction
confidence & 0.50
{[}0.50, 0.55{]} &
0.52 {[}0.50,
0.56{]} & 0.50
{[}0.50, 0.55{]} &
0.51 & wrong
ligand \\
22 & Symmetry swaps
(count) & single
value & Bookkeeping
counter & 0.51
{[}0.50, 0.54{]} &
0.51 {[}0.50,
0.53{]} & 0.50
{[}0.50, 0.52{]} &
0.51 & cognate \\
23 & pLDDT
\cite{jumper2021} &
lowest residue &
Whole-prediction
confidence & 0.50
{[}0.50, 0.52{]} &
0.50 {[}0.50,
0.53{]} & 0.50
{[}0.50, 0.53{]} &
0.50 & cognate \\
24 & AME pass bit
\cite{ahern2026} &
single value &
Active-site accuracy
(pass bit) & 0.50
{[}0.50, 0.50{]} &
0.50 {[}0.50,
0.53{]} & 0.50
{[}0.50, 0.53{]} &
0.50 & mixed \\
25 & Backbone atoms
aligned (count) &
single value &
Bookkeeping counter
& 0.50 {[}0.50,
0.50{]} & 0.50
{[}0.50, 0.50{]} &
0.50 {[}0.50,
0.50{]} & 0.50 &
mixed \\
26 & Catalytic atoms
(count) & single
value & Bookkeeping
counter & 0.50
{[}0.50, 0.50{]} &
0.50 {[}0.50,
0.50{]} & 0.50
{[}0.50, 0.50{]} &
0.50 & mixed \\
27 & Models
predicted (count) &
single value &
Bookkeeping counter
& 0.50 {[}0.50,
0.50{]} & 0.50
{[}0.50, 0.50{]} &
0.50 {[}0.50,
0.50{]} & 0.50 &
mixed \\
28 & Inter-chain
clash flag
\cite{chai2024} &
maximum &
Whole-prediction
confidence & 0.50
{[}0.50, 0.50{]} &
0.50 {[}0.50,
0.50{]} & 0.50
{[}0.50, 0.50{]} &
0.50 & mixed \\
29 & Inter-chain
clash flag
\cite{chai2024} &
mean &
Whole-prediction
confidence & 0.50
{[}0.50, 0.50{]} &
0.50 {[}0.50,
0.50{]} & 0.50
{[}0.50, 0.50{]} &
0.50 & mixed \\
30 & Inter-chain
clash flag
\cite{chai2024} &
spread (SD) &
Whole-prediction
confidence & 0.50
{[}0.50, 0.50{]} &
0.50 {[}0.50,
0.50{]} & 0.50
{[}0.50, 0.50{]} &
0.50 & mixed \\
\end{longtable}

\textbf{More
detail.} The
analysis on all 59
natural enzymes
(five-seed cognate
mean) and on the 30
de novo designs, for
all 33 scored
prediction-based
quantities including
the Chai-1 combined
score, with the
paired change in
units of seed noise:
Table S14. The same
comparison
stratified by the
difference in ligand
size and charge,
with paired win
fractions: Table
S15. The drift of
the current Chai-1
engine against the
earlier predictions
on a re-run shard:
Table S45.

\begin{center}\rule{0.5\linewidth}{0.5pt}\end{center}

\subsection{3.3
Activity ranking
test}

\textbf{Experiment.}
Metrics are scored
on real proteins
whose activity was
measured, and each
metric is ranked by
the rank correlation
between its value
and the measured
activity (Methods
2.4). Two datasets
are used, each with
its own table.

\subsubsection{3.3.1
Ranking natural
enzyme activity}

\textbf{Dataset.}
BglB, a
β-glucosidase from
\emph{Paenibacillus
polymyxa}, with
single-point
variants
characterised by the
D2DCure consortium
\cite{carlin2016}
(data retrieved from
d2dcure.com on
2026-08-06): 655
variants at 207
positions, of which
the \textbf{432
variants at 175
positions that are
folded (expressed
and not
destabilised) and
have kinetics} are
the evaluation set.
The measured value
is the
\emph{impairment},
the log10 of the
loss of kcat/KM
relative to the wild
type (larger = more
impaired). Each
variant is built in
the computer on the
crystal structure of
the wild type (RCSB
PDB entry 2JIE
\cite{burley2025,berman2000},
which carries a
covalent glucosyl
intermediate) and
structure-space
metrics are computed
over the two
catalytic glutamates
(Glu167 and Glu356,
from the UniProt
\cite{uniprot2023}
active-site
features) before and
after the
substitution;
prediction-based
metrics are computed
on Chai-1
\cite{chai2024}
predictions of the
variant with the
assay substrate
(pNPG) (Methods
2.4.1).

\textbf{Ranking
criteria.} The 41
items are the 28
structure-space
main-set metrics
that BglB defines
(the omega term does
not change under a
side-chain
substitution on a
fixed backbone), the
4 prediction-based
main-set metrics, 4
reference rows
(tyrosine fraction,
net charge, pLDDT,
pTM) and 5 baselines
declared before any
metric was joined to
a label: the
distance from the
mutated residue to
the nearest
catalytic residue,
the burial of the
mutated residue, the
BLOSUM62 score of
the substitution,
the change in
side-chain volume
and the Rosetta
score that comes
with the BglB data
(248 of the 432
variants). Each is
ranked by the
Spearman rank
correlation ρ
\cite{spearman1904}
with the measured
impairment, with a
95\% bootstrap
interval that
resamples the 175
positions
\cite{field2007}.
For a metric the
score is the size of
its change (the
absolute difference
between the variant
and the
unsubstituted
protein), with the
direction fixed in
advance (larger in
more impaired
variants); for a
baseline it is the
baseline itself,
pointed so that a
larger value means
more impaired. A ρ
in bold exceeds the
best-of-41 threshold
\cite{westfall1993},
the 95th percentile
of the best ρ among
41 items when the
measured impairments
are permuted among
whole positions
(median 0.24, 95th
percentile 0.32).
The last column says
whether the item
beats the distance
baseline: the
interval of the
difference between
its ρ and the
baseline's, on the
same bootstrap
draws, lies above 0.

\begin{longtable}[]{@{}
  >{\raggedright\arraybackslash}p{(\columnwidth - 10\tabcolsep) * \real{0.0714}}
  >{\raggedright\arraybackslash}p{(\columnwidth - 10\tabcolsep) * \real{0.2937}}
  >{\raggedright\arraybackslash}p{(\columnwidth - 10\tabcolsep) * \real{0.1825}}
  >{\raggedright\arraybackslash}p{(\columnwidth - 10\tabcolsep) * \real{0.1190}}
  >{\raggedright\arraybackslash}p{(\columnwidth - 10\tabcolsep) * \real{0.2063}}
  >{\raggedright\arraybackslash}p{(\columnwidth - 10\tabcolsep) * \real{0.1270}}@{}}
\caption*{\textbf{Table
8. BglB variants:
top 30 of 41 ranked
items on the
432-variant
evaluation set.}
Spearman rank
correlation ρ
between each item
and the measured
impairment of the
432 variants, with a
95\% interval that
resamples the 175
positions. For a
metric the score is
the size of its
change, for a
baseline its own
score, with the
direction fixed in
advance (a larger
change in more
impaired variants);
a negative ρ is
reported as it is.
The 41 items are the
32 main-set metrics
that BglB defines,
four reference rows
and five declared
baselines; reference
rows and baselines
are in italic. A
value in bold
exceeds the
best-of-41 threshold
(median 0.24, 95th
percentile 0.32;
permutation of the
measured impairments
among whole
positions; a
reference level for
reading the
ranking).
Structure-space rows
have 431 variants
with a value,
prediction-based
rows 432 and the
Rosetta score of the
BglB data 248. Items
31 to 41, and all
114 quantities
analysed, are in
Table S16. Source:
main/T8\_bglb\_variants.csv.}\tabularnewline
\toprule\noalign{}
\begin{minipage}[b]{\linewidth}\raggedright
Rank
\end{minipage} &
\begin{minipage}[b]{\linewidth}\raggedright
Item
\end{minipage} &
\begin{minipage}[b]{\linewidth}\raggedright
Group
\end{minipage} &
\begin{minipage}[b]{\linewidth}\raggedright
Mode
\end{minipage} &
\begin{minipage}[b]{\linewidth}\raggedright
Spearman ρ {[}95\%
interval{]}
\end{minipage} &
\begin{minipage}[b]{\linewidth}\raggedright
Beats the distance
baseline
\end{minipage} \\
\midrule\noalign{}
\endfirsthead
\toprule\noalign{}
\begin{minipage}[b]{\linewidth}\raggedright
Rank
\end{minipage} &
\begin{minipage}[b]{\linewidth}\raggedright
Item
\end{minipage} &
\begin{minipage}[b]{\linewidth}\raggedright
Group
\end{minipage} &
\begin{minipage}[b]{\linewidth}\raggedright
Mode
\end{minipage} &
\begin{minipage}[b]{\linewidth}\raggedright
Spearman ρ {[}95\%
interval{]}
\end{minipage} &
\begin{minipage}[b]{\linewidth}\raggedright
Beats the distance
baseline
\end{minipage} \\
\midrule\noalign{}
\endhead
\bottomrule\noalign{}
\endlastfoot
1 & pKa, range
(PROPKA)
\cite{olsson2011,sondergaard2011}
& pKa & structure &
\textbf{0.48
{[}0.35, 0.59{]}} &
no \\
2 & pKa, min
(PROPKA)
\cite{olsson2011,sondergaard2011}
& pKa & structure &
\textbf{0.46
{[}0.32, 0.57{]}} &
no \\
3 & pKa shift, max
absolute (PROPKA)
\cite{olsson2011,sondergaard2011}
& pKa & structure &
\textbf{0.46
{[}0.32, 0.58{]}} &
no \\
4 & pKa, max
(PROPKA)
\cite{olsson2011,sondergaard2011}
& pKa & structure &
\textbf{0.46
{[}0.32, 0.58{]}} &
no \\
5 & \emph{Distance
to the catalytic
residues} &
Reference row & -- &
\textbf{0.46
{[}0.32, 0.57{]}} &
-- \\
6 & Intra-residue
repulsion (Rosetta,
site)
\cite{alford2017} &
Rosetta energy &
structure &
\textbf{0.45
{[}0.30, 0.57{]}} &
no \\
7 & Relative SASA,
mean & Geometry &
structure &
\textbf{0.45
{[}0.32, 0.57{]}} &
no \\
8 & pKa shift, mean
(PROPKA)
\cite{olsson2011,sondergaard2011}
& pKa & structure &
\textbf{0.45
{[}0.31, 0.57{]}} &
no \\
9 & pKa, mean
(PROPKA)
\cite{olsson2011,sondergaard2011}
& pKa & structure &
\textbf{0.45
{[}0.31, 0.57{]}} &
no \\
10 & Total SASA &
Geometry & structure
& \textbf{0.45
{[}0.31, 0.57{]}} &
no \\
11 & Repulsion
(Rosetta, site)
\cite{alford2017} &
Rosetta energy &
structure &
\textbf{0.45
{[}0.31, 0.57{]}} &
no \\
12 & Pairwise
distance, max &
Geometry & structure
& \textbf{0.45
{[}0.30, 0.56{]}} &
no \\
13 & Attraction
(Rosetta, site)
\cite{alford2017} &
Rosetta energy &
structure &
\textbf{0.44
{[}0.30, 0.56{]}} &
no \\
14 & Rotamer,
Dunbrack (Rosetta,
site)
\cite{alford2017} &
Rosetta energy &
structure &
\textbf{0.44
{[}0.29, 0.56{]}} &
no \\
15 & lk-ball
solvation (Rosetta,
site)
\cite{alford2017} &
Rosetta energy &
structure &
\textbf{0.44
{[}0.31, 0.54{]}} &
no \\
16 & Radius of
gyration & Geometry
& structure &
\textbf{0.43
{[}0.28, 0.56{]}} &
no \\
17 & Net van der
Waals (Rosetta,
site)
\cite{alford2017} &
Rosetta energy &
structure &
\textbf{0.43
{[}0.28, 0.55{]}} &
no \\
18 & Relative SASA,
max & Geometry &
structure &
\textbf{0.43
{[}0.29, 0.54{]}} &
no \\
19 & Pairwise
distance, mean &
Geometry & structure
& \textbf{0.43
{[}0.28, 0.55{]}} &
no \\
20 & Electrostatics
(Rosetta, site)
\cite{alford2017} &
Rosetta energy &
structure &
\textbf{0.43
{[}0.28, 0.55{]}} &
no \\
21 & Site total
(Rosetta)
\cite{alford2017} &
Rosetta energy &
structure &
\textbf{0.42
{[}0.27, 0.55{]}} &
no \\
22 & Worst residue
(Rosetta, site)
\cite{alford2017} &
Rosetta energy &
structure &
\textbf{0.42
{[}0.26, 0.55{]}} &
no \\
23 & Solvation
(Rosetta, site)
\cite{alford2017} &
Rosetta energy &
structure &
\textbf{0.41
{[}0.27, 0.54{]}} &
no \\
24 & H-bond, side
chain (Rosetta,
site)
\cite{alford2017} &
Rosetta energy &
structure &
\textbf{0.35
{[}0.19, 0.49{]}} &
no \\
25 & \emph{Burial of
the mutated residue}
& Reference row & --
& \textbf{0.33
{[}0.21, 0.43{]}} &
-- \\
26 & ipTM (Chai-1)
\cite{chai2024,evans2021}
& Interface
confidence &
prediction &
\textbf{0.33
{[}0.19, 0.45{]}} &
no \\
27 & Per-chain pTM
minimum (Chai-1)
\cite{chai2024,jumper2021}
& Interface
confidence &
prediction & 0.28
{[}0.15, 0.41{]} &
no \\
28 & \emph{pTM}
\cite{jumper2021} &
Reference row &
prediction & 0.26
{[}0.12, 0.38{]} &
no \\
29 & \emph{Rosetta
score in the BglB
data set} &
Reference row & -- &
0.21 {[}0.06,
0.35{]} & -- \\
30 & Fraction buried
& Geometry &
structure & 0.19
{[}0.06, 0.30{]} &
no \\
\end{longtable}

\textbf{More
detail.} All 114
quantities analysed,
ranked, with the
AUROC of the binary
contrast (impaired
against
wild-type-like)
beside each rank
correlation: Table
S16. The 84
structure-space and
comparator metrics
with the signed,
direction-free and
magnitude AUROC and
the comparison with
every baseline:
Table S17. The 28
prediction-based
metrics: Table S19.
The main analysis
and its ten
sensitivity
analyses: Table S18.
The five declared
baselines and the
reference rows in
full: Table S12.

\subsubsection{3.3.2
Ranking de novo
designed enzyme
activity}

\textbf{Dataset.}
The 192-design
campaign of 3.1.2
\cite{kim2026}. The
16 designs with a
reported kcat/KM
have a measured
activity and form
the evaluation set,
so every row has n =
16; they lie in 8 of
the 136 backbone
clusters (Methods
2.4.2).

\textbf{Ranking
criteria.} The 36
items are those of
Table 6. Each is
ranked by the size
of the Spearman
correlation ρ
\cite{spearman1904}
between its value on
the design as
deposited and log10
kcat/KM, with a 90\%
bootstrap interval
over designs. The
sign of ρ is
reported (``higher
value means higher
or lower
activity''), and the
last column gives ρ
with the sequence
length partialled
out of both sides. A
ρ in bold has an
interval that
excludes 0, which
about 3.6 of the 36
items would show for
unrelated metrics.
The analysis is
descriptive (Methods
2.4.2).

\begin{longtable}[]{@{}
  >{\raggedright\arraybackslash}p{(\columnwidth - 12\tabcolsep) * \real{0.0652}}
  >{\raggedright\arraybackslash}p{(\columnwidth - 12\tabcolsep) * \real{0.2464}}
  >{\raggedright\arraybackslash}p{(\columnwidth - 12\tabcolsep) * \real{0.1449}}
  >{\raggedright\arraybackslash}p{(\columnwidth - 12\tabcolsep) * \real{0.1014}}
  >{\raggedright\arraybackslash}p{(\columnwidth - 12\tabcolsep) * \real{0.1884}}
  >{\raggedright\arraybackslash}p{(\columnwidth - 12\tabcolsep) * \real{0.1377}}
  >{\raggedright\arraybackslash}p{(\columnwidth - 12\tabcolsep) * \real{0.1159}}@{}}
\caption*{\textbf{Table
9. Plated designs:
top 30 of 36 ranked
items on the 16
designs that have a
reported kcat/KM.}
Spearman rank
correlation ρ
between each metric
(native value on the
design) and log10
kcat/KM, with a 90\%
interval that
resamples designs;
items are ranked by
the size of ρ. A
value in bold has an
interval that
excludes 0 (about
3.6 of the 36 items
would show this for
unrelated metrics).
The last column
gives ρ with the
sequence length
partialled out of
both sides. The 36
items are 29
structure-space and
2 prediction-based
main-set metrics and
5 reference rows
(italic).
Descriptive
analysis. Items 31
to 36, and all 106
metrics with a
correlation, are in
Table S21; the
hits-per-plate
analysis of the 192
designs is in Table
S20. Source:
main/T9\_plated\_designs.csv.}\tabularnewline
\toprule\noalign{}
\begin{minipage}[b]{\linewidth}\raggedright
Rank
\end{minipage} &
\begin{minipage}[b]{\linewidth}\raggedright
Item
\end{minipage} &
\begin{minipage}[b]{\linewidth}\raggedright
Group
\end{minipage} &
\begin{minipage}[b]{\linewidth}\raggedright
Mode
\end{minipage} &
\begin{minipage}[b]{\linewidth}\raggedright
Spearman ρ {[}90\%
interval{]}
\end{minipage} &
\begin{minipage}[b]{\linewidth}\raggedright
Higher value means
\end{minipage} &
\begin{minipage}[b]{\linewidth}\raggedright
ρ, length partialled
out
\end{minipage} \\
\midrule\noalign{}
\endfirsthead
\toprule\noalign{}
\begin{minipage}[b]{\linewidth}\raggedright
Rank
\end{minipage} &
\begin{minipage}[b]{\linewidth}\raggedright
Item
\end{minipage} &
\begin{minipage}[b]{\linewidth}\raggedright
Group
\end{minipage} &
\begin{minipage}[b]{\linewidth}\raggedright
Mode
\end{minipage} &
\begin{minipage}[b]{\linewidth}\raggedright
Spearman ρ {[}90\%
interval{]}
\end{minipage} &
\begin{minipage}[b]{\linewidth}\raggedright
Higher value means
\end{minipage} &
\begin{minipage}[b]{\linewidth}\raggedright
ρ, length partialled
out
\end{minipage} \\
\midrule\noalign{}
\endhead
\bottomrule\noalign{}
\endlastfoot
1 & Polar contacts &
Geometry & structure
& \textbf{0.67
{[}0.41, 0.85{]}} &
higher activity &
0.52 \\
2 & Net van der
Waals (Rosetta,
site)
\cite{alford2017} &
Rosetta energy &
structure &
\textbf{−0.63
{[}−0.88, −0.24{]}}
& lower activity &
−0.47 \\
3 & \emph{pLDDT}
\cite{jumper2021} &
Reference row &
prediction &
\textbf{0.60
{[}0.24, 0.83{]}} &
higher activity &
0.36 \\
4 & \emph{Sequence
length} & Reference
row & structure &
\textbf{−0.58
{[}−0.82, −0.23{]}}
& lower activity &
-- \\
5 & Solvation
(Rosetta, site)
\cite{alford2017} &
Rosetta energy &
structure &
\textbf{0.57
{[}0.15, 0.85{]}} &
higher activity &
0.54 \\
6 & Amino-acid
propensity (Rosetta,
site)
\cite{alford2017} &
Rosetta energy &
structure &
\textbf{0.54
{[}0.08, 0.81{]}} &
higher activity &
0.64 \\
7 & \emph{Net
charge} & Reference
row & structure &
0.46 {[}−0.02,
0.80{]} & higher
activity & 0.22 \\
8 & Pairwise
distance, mean &
Geometry & structure
& −0.43 {[}−0.71,
0.02{]} & lower
activity & −0.29 \\
9 & Site total
(Rosetta)
\cite{alford2017} &
Rosetta energy &
structure & −0.42
{[}−0.69, 0.01{]} &
lower activity &
−0.37 \\
10 & Radius of
gyration & Geometry
& structure & −0.41
{[}−0.69, 0.03{]} &
lower activity &
−0.29 \\
11 & H-bond, side
chain (Rosetta,
site)
\cite{alford2017} &
Rosetta energy &
structure & −0.40
{[}−0.74, 0.05{]} &
lower activity &
−0.14 \\
12 & Ramachandran
(Rosetta, site)
\cite{alford2017} &
Rosetta energy &
structure & −0.38
{[}−0.82, 0.15{]} &
lower activity &
−0.09 \\
13 & Repulsion
(Rosetta, site)
\cite{alford2017} &
Rosetta energy &
structure & −0.33
{[}−0.73, 0.17{]} &
lower activity &
0.21 \\
14 & pKa, mean
(PROPKA)
\cite{olsson2011,sondergaard2011}
& pKa & structure &
−0.30 {[}−0.57,
0.09{]} & lower
activity & 0.09 \\
15 & pKa shift, mean
(PROPKA)
\cite{olsson2011,sondergaard2011}
& pKa & structure &
−0.30 {[}−0.56,
0.09{]} & lower
activity & 0.09 \\
16 & Relative SASA,
mean & Geometry &
structure & −0.25
{[}−0.76, 0.31{]} &
lower activity &
−0.35 \\
17 & Total SASA &
Geometry & structure
& −0.25 {[}−0.76,
0.30{]} & lower
activity & −0.35 \\
18 & Relative SASA,
max & Geometry &
structure & −0.25
{[}−0.76, 0.28{]} &
lower activity &
−0.34 \\
19 & Attraction
(Rosetta, site)
\cite{alford2017} &
Rosetta energy &
structure & −0.25
{[}−0.66, 0.25{]} &
lower activity &
−0.20 \\
20 & Pairwise
distance, max &
Geometry & structure
& 0.23 {[}−0.26,
0.64{]} & higher
activity & 0.26 \\
21 & Intra-residue
repulsion (Rosetta,
site)
\cite{alford2017} &
Rosetta energy &
structure & −0.23
{[}−0.65, 0.30{]} &
lower activity &
−0.46 \\
22 & lk-ball
solvation (Rosetta,
site)
\cite{alford2017} &
Rosetta energy &
structure & −0.21
{[}−0.53, 0.19{]} &
lower activity &
0.08 \\
23 & Electrostatics
(Rosetta, site)
\cite{alford2017} &
Rosetta energy &
structure & −0.19
{[}−0.58, 0.25{]} &
lower activity &
−0.61 \\
24 & Per-chain pTM
minimum (Chai-1)
\cite{chai2024,jumper2021}
& Interface
confidence &
prediction & 0.18
{[}−0.35, 0.66{]} &
higher activity &
−0.00 \\
25 & Rotamer,
Dunbrack (Rosetta,
site)
\cite{alford2017} &
Rosetta energy &
structure & 0.18
{[}−0.31, 0.57{]} &
higher activity &
0.11 \\
26 & H-bond,
backbone to side
chain (Rosetta,
site)
\cite{alford2017} &
Rosetta energy &
structure & 0.17
{[}−0.21, 0.52{]} &
higher activity &
0.32 \\
27 & pKa, max
(PROPKA)
\cite{olsson2011,sondergaard2011}
& pKa & structure &
−0.16 {[}−0.49,
0.26{]} & lower
activity & 0.16 \\
28 & \emph{pTM}
\cite{jumper2021} &
Reference row &
prediction & −0.15
{[}−0.64, 0.41{]} &
lower activity &
−0.26 \\
29 & \emph{Tyrosine
fraction} &
Reference row &
structure & −0.14
{[}−0.68, 0.40{]} &
lower activity &
0.00 \\
30 & pKa, range
(PROPKA)
\cite{olsson2011,sondergaard2011}
& pKa & structure &
−0.13 {[}−0.53,
0.31{]} & lower
activity & −0.07 \\
\end{longtable}

\textbf{More
detail.} The
correlation with
kcat/KM of all 106
metrics that have
one, with the length
control: Table S21.
The correlations of
the structure-space
metrics and
comparators, with
and without the
length control:
Table S22. The hits
per plate of all 128
metrics when each
fills a 96-well
plate from the
quarter of the 192
designs that it
ranks highest or
lowest: Table S20.
The enrichment at
keep fractions of 5,
10, 25 and 50\% in
both directions:
Table S23 (102
structure-space
metrics) and Table
S24 (29
prediction-based
metrics).

\begin{center}\rule{0.5\linewidth}{0.5pt}\end{center}

\subsection{3.4
Detection ability
test}

\textbf{Experiment.}
Each metric is
ranked by whether it
changes more under a
lesion at a
catalytic residue
than under the
identical lesion at
a matched
non-catalytic
control position of
the same protein
(Methods 2.5). Two
kinds of lesion are
tested, each with
its own datasets and
table: the catalytic
lesion (3.4.1) and
the second-shell
lesion (3.4.2).

\subsubsection{3.4.1
Detecting catalytic
lesion}

\textbf{Experiment.}
One catalytic
residue is replaced
in four steps of
increasing change in
side-chain volume
(isosteric 15 ų,
non-isosteric 36 ų,
to Ala 53 ų, to Gly
81 ų), in the
crystal structure by
a rotamer-space edit
and repack, or in
the sequence of the
re-predicted complex
with Chai-1
\cite{chai2024}; the
control makes the
identical
substitution at a
matched buried
non-catalytic
position of the same
protein.

\textbf{Dataset.}
For the 29
structure-space
main-set metrics:
the 143 M-CSA
enzymes
\cite{ribeiro2018}
(77 enzyme
sub-subclasses;
packing-matched,
count-gated
controls) pooled
with the 52 enzymes
of the 53-enzyme
pilot set that are
not in it (active
sites from UniProt
\cite{uniprot2023}
and M-CSA; controls
matched on residue
type and burial
only; the one pilot
enzyme that shares a
PDB entry with the
main set is counted
once), 195 enzymes
in all. For the four
prediction-based
main-set metrics and
the two reference
rows (pLDDT and
pTM): the natural
arm of the
re-predicted ladder,
the 59 natural
enzymes of the
substrate-swap set
with 286
catalytic/control
pairs. The 30 de
novo designs and the
15 designed systems
of the ladder, whose
control is
constructed, are
analysed separately
(Tables S29 and S30)
(Methods 2.5.1).

\textbf{Ranking
criteria.} The 35
items are the 29
structure-space and
4 prediction-based
main-set metrics and
the 2 reference rows
(italic). The rank
statistic R is the
probability that a
metric changes more
at the catalytic
lesion than at the
matched control of
the same protein
(0.5 = no
specificity; above
0.5, more change at
the catalytic
lesion). The change
of a metric is the
absolute difference
between the
perturbed and the
unperturbed
structure (or
prediction), a tie
counts one half, and
proteins in which
neither arm moves
are dropped. A step
counts toward a
structure-space
metric only if the
metric responds to
the lesion (the
interval of its
median change
excludes zero) and
the control response
is above the
metric's own noise;
a step that is not
counted is in
parentheses, and
``--'' is a step
with fewer than five
informative
proteins. The mean R
over the counted
steps has a 95\%
bootstrap interval
that resamples
enzyme
sub-subclasses
(systems for the
predictor). Rows are
grouped by the
lesion at which the
metric is first
detected (the
smallest step at
which it responds;
``Never'' if at
none) and ranked by
mean R within each
group; a metric with
no counted step has
no mean and closes
its group. A mean in
bold exceeds the
best-of-33 threshold
\cite{westfall1993},
the 95th percentile
of the best R among
33 items when the
catalytic and
control labels are
swapped within each
protein (median
0.57, 95th
percentile 0.62).
The distance-matched
R (prediction-based
metrics only) is
computed on the
pairs whose control
lies at the same
distance from the
ligand as the
catalytic residue
(within 2 Å and
within 0.10 in
relative
solvent-accessible
area), pooled over
the four steps (68
to 69 pairs in 21 to
22 systems); the
original control of
the predictor sat a
median 13.4 Å from
the ligand, the
catalytic residue
3.1 Å (the distances
for every pair are
in Table S37 and the
caliper-matched
subsets in Table
S40). The number of
steps counted for
each metric and the
smallest step at
which a
structure-space
metric is specific
(it responds, the
specificity ratio is
above 1, the control
response is above
the noise and both
arms changed equal
numbers of residues)
are in Table S28.

\begin{longtable}[]{@{}
  >{\raggedright\arraybackslash}p{(\columnwidth - 10\tabcolsep) * \real{0.1007}}
  >{\raggedright\arraybackslash}p{(\columnwidth - 10\tabcolsep) * \real{0.0647}}
  >{\raggedright\arraybackslash}p{(\columnwidth - 10\tabcolsep) * \real{0.2590}}
  >{\raggedright\arraybackslash}p{(\columnwidth - 10\tabcolsep) * \real{0.2158}}
  >{\raggedright\arraybackslash}p{(\columnwidth - 10\tabcolsep) * \real{0.1727}}
  >{\raggedright\arraybackslash}p{(\columnwidth - 10\tabcolsep) * \real{0.1871}}@{}}
\caption*{\textbf{Table
10. Catalytic lesion
(chemistry ladder):
the 35 ranked items
grouped by the
lesion at which they
are first detected
and ranked by R
within the group.}
Rank statistic R,
the probability that
a metric changes
more at the
catalytic lesion
than at the matched
control (0.5 = no
specificity), at the
four steps and as
the mean over the
steps counted, with
a 95\% interval that
resamples enzyme
sub-subclasses
(systems for the
predictor).
``Detected from'' is
the smallest lesion
step (isosteric 15
ų, non-isosteric 36
ų, Ala 53 ų, Gly
81 ų) at which the
metric responds.
Structure-space
metrics: up to 195
natural enzymes (143
main + 52 pilot; the
number of enzymes in
which a metric moves
at all ranges from 1
to 150 per metric
and step);
prediction-based
metrics (ipTM,
per-chain pTM
minimum, ligand
clearance, AME RMSD)
and the reference
rows pLDDT and pTM
(italic): 59 natural
enzymes, 286 pairs.
A step in
parentheses is not
counted. A mean in
bold exceeds the
best-of-33 threshold
(median 0.57, 95th
percentile 0.62; a
reference level for
reading the
ranking). The
distance-matched R
is given for the
prediction-based
metrics. The 8
supplementary
comparators, the
number of steps
counted, the onset
of specificity and
all per-step flags
are in Table S28.
Source:
main/T10\_specificity\_catalytic\_lesion.csv.}\tabularnewline
\toprule\noalign{}
\begin{minipage}[b]{\linewidth}\raggedright
Detected from
\end{minipage} &
\begin{minipage}[b]{\linewidth}\raggedright
Rank in group
\end{minipage} &
\begin{minipage}[b]{\linewidth}\raggedright
Item
\end{minipage} &
\begin{minipage}[b]{\linewidth}\raggedright
R at isosteric /
non-isosteric / Ala
/ Gly
\end{minipage} &
\begin{minipage}[b]{\linewidth}\raggedright
Mean R {[}95\%
interval{]}
\end{minipage} &
\begin{minipage}[b]{\linewidth}\raggedright
R, distance-matched
{[}95\% interval{]}
\end{minipage} \\
\midrule\noalign{}
\endfirsthead
\toprule\noalign{}
\begin{minipage}[b]{\linewidth}\raggedright
Detected from
\end{minipage} &
\begin{minipage}[b]{\linewidth}\raggedright
Rank in group
\end{minipage} &
\begin{minipage}[b]{\linewidth}\raggedright
Item
\end{minipage} &
\begin{minipage}[b]{\linewidth}\raggedright
R at isosteric /
non-isosteric / Ala
/ Gly
\end{minipage} &
\begin{minipage}[b]{\linewidth}\raggedright
Mean R {[}95\%
interval{]}
\end{minipage} &
\begin{minipage}[b]{\linewidth}\raggedright
R, distance-matched
{[}95\% interval{]}
\end{minipage} \\
\midrule\noalign{}
\endhead
\bottomrule\noalign{}
\endlastfoot
Isosteric & 1 &
Per-chain pTM
minimum (Chai-1)
\cite{chai2024,jumper2021}
& 0.81 / 0.73 / 0.75
/ 0.72 &
\textbf{0.75
{[}0.68, 0.82{]}} &
0.59 {[}0.42,
0.76{]} \\
Isosteric & 2 & pKa,
range (PROPKA)
\cite{olsson2011,sondergaard2011}
& 0.70 / 0.67 / -- /
-- & \textbf{0.68
{[}0.58, 0.78{]}} &
-- \\
Isosteric & 3 & ipTM
(Chai-1)
\cite{chai2024,evans2021}
& 0.71 / 0.71 / 0.69
/ 0.58 &
\textbf{0.67
{[}0.60, 0.75{]}} &
0.35 {[}0.24,
0.50{]} \\
Isosteric & 4 &
Solvation (Rosetta,
site)
\cite{alford2017} &
0.59 / 0.73 / 0.60 /
0.62 & \textbf{0.63
{[}0.57, 0.70{]}} &
-- \\
Isosteric & 5 &
Relative SASA, mean
& 0.48 / 0.55 / 0.73
/ 0.69 & 0.61
{[}0.56, 0.67{]} &
-- \\
Isosteric & 6 & pKa,
min (PROPKA)
\cite{olsson2011,sondergaard2011}
& 0.54 / 0.67 / -- /
-- & 0.61 {[}0.53,
0.67{]} & -- \\
Isosteric & 7 &
Repulsion (Rosetta,
site)
\cite{alford2017} &
0.48 / 0.55 / 0.64 /
0.62 & 0.57 {[}0.52,
0.62{]} & -- \\
Isosteric & 8 & pKa
shift, mean (PROPKA)
\cite{olsson2011,sondergaard2011}
& 0.46 / 0.67 / -- /
-- & 0.56 {[}0.47,
0.66{]} & -- \\
Isosteric & 9 &
\emph{pTM}
\cite{jumper2021} &
0.60 / 0.48 / 0.60 /
0.58 & 0.56 {[}0.49,
0.64{]} & 0.45
{[}0.30, 0.60{]} \\
Isosteric & 10 &
Relative SASA, max &
0.49 / (0.50) / 0.56
/ 0.55 & 0.53
{[}0.48, 0.59{]} &
-- \\
Isosteric & 11 &
Attraction (Rosetta,
site)
\cite{alford2017} &
0.53 / 0.58 / 0.55 /
0.46 & 0.53 {[}0.46,
0.60{]} & -- \\
Isosteric & 12 &
\emph{pLDDT}
\cite{jumper2021} &
0.62 / 0.37 / 0.55 /
0.58 & 0.53 {[}0.46,
0.60{]} & 0.49
{[}0.34, 0.63{]} \\
Isosteric & 13 &
Ramachandran
(Rosetta, site)
\cite{alford2017} &
0.43 / 0.55 / 0.55 /
0.56 & 0.52 {[}0.47,
0.59{]} & -- \\
Isosteric & 14 &
Intra-residue
repulsion (Rosetta,
site)
\cite{alford2017} &
0.48 / 0.59 / 0.48 /
0.52 & 0.52 {[}0.47,
0.57{]} & -- \\
Isosteric & 15 &
Rotamer, Dunbrack
(Rosetta, site)
\cite{alford2017} &
0.45 / 0.47 / 0.58 /
0.58 & 0.52 {[}0.46,
0.57{]} & -- \\
Isosteric & 16 &
Radius of gyration &
0.42 / 0.48 / 0.62 /
-- & 0.51 {[}0.45,
0.57{]} & -- \\
Isosteric & 17 &
Amino-acid
propensity (Rosetta,
site)
\cite{alford2017} &
0.45 / (0.57) / 0.52
/ 0.53 & 0.50
{[}0.45, 0.55{]} &
-- \\
Isosteric & 18 &
pKa, max (PROPKA)
\cite{olsson2011,sondergaard2011}
& 0.43 / 0.56 / -- /
-- & 0.49 {[}0.39,
0.61{]} & -- \\
Isosteric & 19 &
Worst residue
(Rosetta, site)
\cite{alford2017} &
0.45 / 0.46 / 0.54 /
(0.43) & 0.48
{[}0.44, 0.53{]} &
-- \\
Isosteric & 20 & Net
van der Waals
(Rosetta, site)
\cite{alford2017} &
0.45 / 0.55 / 0.47 /
0.39 & 0.46 {[}0.39,
0.52{]} & -- \\
Isosteric & 21 &
H-bond, side chain
(Rosetta, site)
\cite{alford2017} &
0.47 / 0.42 / 0.43 /
0.42 & 0.43 {[}0.35,
0.52{]} & -- \\
Isosteric & 22 &
Site total (Rosetta)
\cite{alford2017} &
0.41 / 0.58 / 0.41 /
0.29 & 0.42 {[}0.37,
0.47{]} & -- \\
Isosteric & 23 &
Pairwise distance,
mean & 0.43 / 0.49 /
0.32 / -- & 0.41
{[}0.34, 0.50{]} &
-- \\
Isosteric & 24 &
pKa, mean (PROPKA)
\cite{olsson2011,sondergaard2011}
& 0.40 / (0.47) / --
/ -- & 0.40 {[}0.28,
0.54{]} & -- \\
Isosteric & 25 &
H-bond, backbone to
side chain (Rosetta,
site)
\cite{alford2017} &
0.26 / 0.29 / 0.31 /
0.29 & 0.29 {[}0.19,
0.39{]} & -- \\
Isosteric & 26 &
Electrostatics
(Rosetta, site)
\cite{alford2017} &
0.26 / 0.20 / 0.24 /
0.26 & 0.24 {[}0.18,
0.30{]} & -- \\
Isosteric & -- &
Polar contacts &
(0.93) / (0.97) /
(0.96) / (0.96) & --
& -- \\
Non-isosteric & 1 &
AME RMSD
\cite{ahern2026} &
0.61 / 0.73 / 0.68 /
0.72 & \textbf{0.68
{[}0.62, 0.75{]}} &
0.45 {[}0.33,
0.57{]} \\
Non-isosteric & 2 &
lk-ball solvation
(Rosetta, site)
\cite{alford2017} &
(0.56) / 0.53 / 0.59
/ 0.65 & 0.59
{[}0.51, 0.67{]} &
-- \\
Non-isosteric & 3 &
Total SASA & (0.46)
/ 0.51 / (0.49) /
0.58 & 0.55 {[}0.48,
0.60{]} & -- \\
Non-isosteric & 4 &
pKa shift, max
absolute (PROPKA)
\cite{olsson2011,sondergaard2011}
& (0.53) / 0.53 / --
/ -- & 0.53 {[}0.42,
0.65{]} & -- \\
Ala & 1 & Fraction
buried & (0.54) /
(0.56) / 0.70 / 0.63
& \textbf{0.66
{[}0.59, 0.74{]}} &
-- \\
Gly & 1 & Omega
torsion (Rosetta,
site)
\cite{alford2017} &
(0.46) / -- / (0.58)
/ 0.52 & 0.52
{[}0.45, 0.60{]} &
-- \\
Never & 1 & Ligand
clearance
\cite{ahern2026} &
0.57 / 0.63 / 0.57 /
0.63 & 0.60 {[}0.54,
0.66{]} & 0.57
{[}0.44, 0.70{]} \\
Never & -- &
Pairwise distance,
max & (0.30) /
(0.36) / (0.19) / --
& -- & -- \\
\end{longtable}

\textbf{More
detail.}

\begin{itemize}
\tightlist
\item
  \emph{Ranking.}
  All 43 metrics
  scored for the
  catalytic lesion,
  with the 8
  supplementary
  comparators, R and
  the response,
  control-noise and
  specificity flags
  at each step, and
  the mean R of the
  143 main enzymes
  alone: Table S28.
  The counts of
  eligible,
  responding and
  specific
  structure-space
  metrics by family
  and step, with the
  panel statistic
  and the four
  sensitivity sets:
  Table S29.
\item
  \emph{Specificity
  ratio.} The ratio,
  its interval and
  the verdict of
  every eligible
  metric at every
  step: Table S33
  (comparators:
  Table S34); the 97
  cells that survive
  the controls:
  Table S35; R per
  step on the 143
  main enzymes with
  90\% intervals:
  Table S36; the
  robustness to
  motif size,
  packing matching
  and count gating:
  Table S31; the
  panel statistic
  against the
  earlier
  whole-panel
  analysis: Table
  S42; the effective
  number of
  independent
  metrics: Table
  S46.
\item
  \emph{Controls and
  floors.} The
  loosest matching
  tier needed, and
  the balance of
  residues changed
  in the two arms:
  Tables S26 and
  S27; the burial of
  the constructed
  control in the de
  novo designs:
  Table S47; the
  sanity floors:
  Table S25.
\item
  \emph{Predictor
  ladder.} Every
  metric at every
  step: Table S43;
  the three controls
  and the
  pre-specified
  rule: Table S30;
  the distance of
  each pair to the
  ligand: Table S37;
  the ratio in
  distance strata:
  Table S39; the
  caliper-matched
  subsets: Table
  S40; the
  regression
  adjustment: Table
  S41; the 54 pairs
  given a new
  distance-matched
  control: Table
  S38; the seed
  noise: Table S32.
\item
  \emph{Not made.}
  Experiments that
  could not be made
  or were retired:
  Table S48; PLACER
  on the isosteric
  step: Table S49.
\end{itemize}

\subsubsection{3.4.2
Detecting second
shell lesion}

\textbf{Experiment.}
One, two or four
residues of the
second shell around
the catalytic
residues are mutated
and the structure is
repacked
(second-shell
removal); the
control does the
same for the same
number of residues
around matched
non-catalytic
positions, and both
arms are scored over
the catalytic
residues.

\textbf{Dataset.}
The 143 M-CSA
enzymes
\cite{ribeiro2018}
pooled with the 52
pilot enzymes that
are not in it, with
a control arm for up
to 193 of them; the
number of enzymes in
which a metric moves
at all ranges from 0
to 182 per metric
and dose.
Structure-space
metrics are scored
on this lesion
(Methods 2.5.2).

\textbf{Ranking
criteria.} The 29
structure-space
main-set metrics are
ranked as in Table
10, with R at the
three doses (1, 2
and 4 residues
removed) and its
mean over the doses;
``Detected from'' is
the smallest number
of residues removed
at which the metric
responds (``Never''
if none), and rows
within a group are
ranked by raw R. A
metric with no R in
at least five
proteins has ``--''.
Raw values are
reported (see the
caption); the last
column counts, for
each metric, the
doses (of 3) at
which it responds to
the removal and at
which the control
response is above
the metric's noise.

\begin{longtable}[]{@{}
  >{\raggedright\arraybackslash}p{(\columnwidth - 12\tabcolsep) * \real{0.0851}}
  >{\raggedright\arraybackslash}p{(\columnwidth - 12\tabcolsep) * \real{0.0638}}
  >{\raggedright\arraybackslash}p{(\columnwidth - 12\tabcolsep) * \real{0.2128}}
  >{\raggedright\arraybackslash}p{(\columnwidth - 12\tabcolsep) * \real{0.1277}}
  >{\raggedright\arraybackslash}p{(\columnwidth - 12\tabcolsep) * \real{0.1915}}
  >{\raggedright\arraybackslash}p{(\columnwidth - 12\tabcolsep) * \real{0.1631}}
  >{\raggedright\arraybackslash}p{(\columnwidth - 12\tabcolsep) * \real{0.1560}}@{}}
\caption*{\textbf{Table
11. Second-shell
lesion: the 29
structure-space
main-set metrics
grouped by the dose
at which they are
first detected and
ranked by raw R
within the group.}
Rank statistic R,
the probability that
a metric changes
more when
second-shell
residues are removed
around the catalytic
residues than around
a matched control
(0.5 = no
specificity), at the
three doses and as
the mean over the
doses, with a 95\%
interval that
resamples enzyme
sub-subclasses.
``Detected from'' is
the smallest number
of residues removed
at which the metric
responds (``Never''
if none); a metric
with no R in at
least five proteins
has ``--''. No cell
qualifies as
interpretable, so
raw values are
reported: the last
column counts, for
each metric, the
doses (of 3) at
which it responds to
the removal and at
which the control
response is above
the metric's noise.
The 8 supplementary
comparators and the
per-dose flags are
in Table S44.
Source:
main/T11\_specificity\_second\_shell\_lesion.csv.}\tabularnewline
\toprule\noalign{}
\begin{minipage}[b]{\linewidth}\raggedright
Detected from
\end{minipage} &
\begin{minipage}[b]{\linewidth}\raggedright
Rank in group
\end{minipage} &
\begin{minipage}[b]{\linewidth}\raggedright
Item
\end{minipage} &
\begin{minipage}[b]{\linewidth}\raggedright
Group
\end{minipage} &
\begin{minipage}[b]{\linewidth}\raggedright
R at 1 / 2 / 4
residues removed
\end{minipage} &
\begin{minipage}[b]{\linewidth}\raggedright
Mean R {[}95\%
interval{]}
\end{minipage} &
\begin{minipage}[b]{\linewidth}\raggedright
Doses with a
response / with a
control above noise
\end{minipage} \\
\midrule\noalign{}
\endfirsthead
\toprule\noalign{}
\begin{minipage}[b]{\linewidth}\raggedright
Detected from
\end{minipage} &
\begin{minipage}[b]{\linewidth}\raggedright
Rank in group
\end{minipage} &
\begin{minipage}[b]{\linewidth}\raggedright
Item
\end{minipage} &
\begin{minipage}[b]{\linewidth}\raggedright
Group
\end{minipage} &
\begin{minipage}[b]{\linewidth}\raggedright
R at 1 / 2 / 4
residues removed
\end{minipage} &
\begin{minipage}[b]{\linewidth}\raggedright
Mean R {[}95\%
interval{]}
\end{minipage} &
\begin{minipage}[b]{\linewidth}\raggedright
Doses with a
response / with a
control above noise
\end{minipage} \\
\midrule\noalign{}
\endhead
\bottomrule\noalign{}
\endlastfoot
4 residues & 1 & pKa
shift, mean (PROPKA)
\cite{olsson2011,sondergaard2011}
& pKa & 0.74 / 0.72
/ 0.64 & 0.70
{[}0.64, 0.75{]} & 1
/ 0 \\
4 residues & 2 &
pKa, mean (PROPKA)
\cite{olsson2011,sondergaard2011}
& pKa & 0.73 / 0.71
/ 0.64 & 0.69
{[}0.63, 0.75{]} & 1
/ 0 \\
Never & 1 & Fraction
buried & Geometry &
0.91 / 0.79 / 0.65 &
0.78 {[}0.68,
0.87{]} & 0 / 0 \\
Never & 2 & Rotamer,
Dunbrack (Rosetta,
site)
\cite{alford2017} &
Rosetta energy &
0.79 / 0.76 / 0.69 &
0.75 {[}0.69,
0.79{]} & 0 / 0 \\
Never & 3 & pKa, max
(PROPKA)
\cite{olsson2011,sondergaard2011}
& pKa & 0.75 / 0.72
/ 0.68 & 0.71
{[}0.66, 0.77{]} & 0
/ 0 \\
Never & 4 & Pairwise
distance, mean &
Geometry & 0.76 /
0.75 / 0.63 & 0.71
{[}0.65, 0.77{]} & 0
/ 0 \\
Never & 5 & Radius
of gyration &
Geometry & 0.77 /
0.71 / 0.64 & 0.71
{[}0.64, 0.76{]} & 0
/ 0 \\
Never & 6 & pKa,
range (PROPKA)
\cite{olsson2011,sondergaard2011}
& pKa & 0.74 / 0.72
/ 0.65 & 0.71
{[}0.65, 0.76{]} & 0
/ 0 \\
Never & 7 & lk-ball
solvation (Rosetta,
site)
\cite{alford2017} &
Rosetta energy &
0.72 / 0.73 / 0.66 &
0.70 {[}0.66,
0.75{]} & 0 / 1 \\
Never & 8 & Total
SASA & Geometry &
0.73 / 0.71 / 0.66 &
0.70 {[}0.65,
0.75{]} & 0 / 0 \\
Never & 9 &
Intra-residue
repulsion (Rosetta,
site)
\cite{alford2017} &
Rosetta energy &
0.72 / 0.73 / 0.64 &
0.70 {[}0.64,
0.75{]} & 0 / 0 \\
Never & 10 & pKa
shift, max absolute
(PROPKA)
\cite{olsson2011,sondergaard2011}
& pKa & 0.75 / 0.68
/ 0.66 & 0.70
{[}0.64, 0.75{]} & 0
/ 1 \\
Never & 11 &
Relative SASA, mean
& Geometry & 0.73 /
0.71 / 0.64 & 0.70
{[}0.64, 0.74{]} & 0
/ 0 \\
Never & 12 &
Attraction (Rosetta,
site)
\cite{alford2017} &
Rosetta energy &
0.73 / 0.69 / 0.61 &
0.67 {[}0.63,
0.72{]} & 0 / 0 \\
Never & 13 &
Relative SASA, max &
Geometry & 0.69 /
0.69 / 0.64 & 0.67
{[}0.61, 0.73{]} & 0
/ 0 \\
Never & 14 &
Electrostatics
(Rosetta, site)
\cite{alford2017} &
Rosetta energy &
0.74 / 0.70 / 0.58 &
0.67 {[}0.62,
0.72{]} & 0 / 1 \\
Never & 15 & Worst
residue (Rosetta,
site)
\cite{alford2017} &
Rosetta energy &
0.74 / 0.68 / 0.59 &
0.67 {[}0.61,
0.72{]} & 0 / 0 \\
Never & 16 & pKa,
min (PROPKA)
\cite{olsson2011,sondergaard2011}
& pKa & 0.74 / 0.67
/ 0.60 & 0.67
{[}0.61, 0.73{]} & 0
/ 0 \\
Never & 17 & Net van
der Waals (Rosetta,
site)
\cite{alford2017} &
Rosetta energy &
0.74 / 0.67 / 0.58 &
0.66 {[}0.61,
0.71{]} & 0 / 1 \\
Never & 18 & H-bond,
side chain (Rosetta,
site)
\cite{alford2017} &
Rosetta energy &
0.72 / 0.68 / 0.59 &
0.66 {[}0.59,
0.72{]} & 0 / 1 \\
Never & 19 & Site
total (Rosetta)
\cite{alford2017} &
Rosetta energy &
0.71 / 0.66 / 0.60 &
0.66 {[}0.60,
0.71{]} & 0 / 1 \\
Never & 20 &
Pairwise distance,
max & Geometry &
0.72 / 0.65 / 0.60 &
0.66 {[}0.52,
0.76{]} & 0 / 0 \\
Never & 21 &
Repulsion (Rosetta,
site)
\cite{alford2017} &
Rosetta energy &
0.71 / 0.63 / 0.62 &
0.65 {[}0.60,
0.71{]} & 0 / 1 \\
Never & 22 &
Solvation (Rosetta,
site)
\cite{alford2017} &
Rosetta energy &
0.72 / 0.66 / 0.57 &
0.65 {[}0.60,
0.70{]} & 0 / 1 \\
Never & 23 & H-bond,
backbone to side
chain (Rosetta,
site)
\cite{alford2017} &
Rosetta energy &
0.56 / 0.68 / 0.65 &
0.63 {[}0.56,
0.70{]} & 0 / 0 \\
Never & 24 & Polar
contacts & Geometry
& 0.59 / 0.59 / 0.65
& 0.61 {[}0.49,
0.73{]} & 0 / 0 \\
Never & 25 &
Ramachandran
(Rosetta, site)
\cite{alford2017} &
Rosetta energy &
0.10 / 0.00 / 0.00 &
0.03 {[}0.00,
0.11{]} & 0 / 3 \\
Never & -- & Omega
torsion (Rosetta,
site)
\cite{alford2017} &
Rosetta energy & --
/ -- / -- & -- & 0 /
3 \\
Never & -- &
Amino-acid
propensity (Rosetta,
site)
\cite{alford2017} &
Rosetta energy & --
/ -- / -- & -- & 0 /
3 \\
\end{longtable}

\textbf{More
detail.} The 37
metrics scored for
the second-shell
lesion, with the 8
supplementary
comparators and the
response and
control-noise flags
at each dose: Table
S44. The specificity
ratio and verdict at
each dose: Table S33
(comparators: Table
S34). The counts of
eligible, responding
and specific metrics
by family and dose:
Table S29.

\begin{center}\rule{0.5\linewidth}{0.5pt}\end{center}

\section{4
Discussion}

\bibliographystyle{unsrt}
\bibliography{references}

\end{document}
